\documentclass[acmsmall,screen,nonacm]{acmart}

\usepackage{booktabs}
\usepackage{microtype}
\usepackage{listings}
\lstset{basicstyle=\ttfamily\small,columns=fullflexible,keepspaces=true,frame=none,commentstyle=\upshape}

\title{Compiling Against Intent}
\author{Or Magen}
\affiliation{\institution{Independent Researcher}\country{USA}}

\begin{abstract}
Source programs mix invariant requirements with incidental choices, and an
optimizing compiler bound by the as-if rule must preserve both. The set of
programs that satisfy the requirements is often larger than the set of as-if
variants of a correct baseline, so as-if search may omit a cheaper legal
program. A queue whose requirements are FIFO order and capacity at least 1000
may be written with 1000 cells and remainder indexing. A 1024-cell ring indexed
by a mask still satisfies those requirements, and as-if cannot emit that ring
from the baseline. That set cannot be searched as a whole.

\textsc{IntentCompile} is a restricted compiler that starts from an existing
baseline and a checkable contract naming those requirements, generates
candidates, admits those a sound verifier proves, and ranks them by
measurement. \textsc{IntC} is a C instance: a rewrite catalog, CBMC on bounded
harnesses, and MSVC \texttt{/O2}. On that baseline, \texttt{/O2} keeps 1000
cells and remainder indexing, and \textsc{IntC} admits the 1024-cell program
and selects it at \(2.3\times\) on the named benchmark.
\end{abstract}

\theoremstyle{acmdefinition}
\newtheorem{algo}{Algorithm}

\begin{document}
\maketitle

\section{Introduction}
\label{sec:intro}

An optimizing compiler may replace a program with any program that has the same
observable behavior, commonly referred to as the as-if rule of ISO
C++~\cite{n4950}. That rule lets compilers reorder, inline, vectorize,
and eliminate code, and it makes the source the complete specification. Every
detail the source fixes is one the compiler must preserve, whether the developer
cared about it or not.

Source code mixes two kinds of decisions. Domain invariants are what the program
must actually guarantee, such as packets leaving a queue in the order they
arrived and the queue holding at least a thousand packets before it drops.
Incidental choices are details the code had to fix in order to run, such as a
capacity of exactly 1000 or a remainder instruction that maps a counter onto
that capacity. A packet queue written for bursts of 800 packets with 25\%
headroom needs at least 1000 slots. The developer writes 1000, and each enqueue
computes \texttt{slot = i \% 1000}, which compiles to a chain of multiplications
and shifts. With 1024 slots the same index is a bitwise \texttt{and}. Both sizes
meet the requirement, and the as-if rule allows only the size in the source,
because the two queues disagree once they hold 1000 packets.

Developers split these decisions by hand. Folly's single-producer queue rounds
its backing array up to a power of two and indexes with a mask, while
\texttt{isFull} still respects the capacity the caller asked
for~\cite{folly_pcq}. DPDK's exact-size rings do the same
separation~\cite{dpdk_rte_ring}. ClickHouse sizes hash tables in powers of two
and places keys with a mask~\cite{clickhouse_hashtable}. Those libraries write
the split into the source, and a compiler bound by as-if cannot introduce it from
a baseline that used 1000 and a remainder. This paper asks for that compiler.

Compiling against a specification instead of the source is an old idea. Leroy
defines compilers that preserve a specification rather than exact
behavior~\cite{leroy2009backend}. Refinement calculi derive implementations from
specifications~\cite{back1998refinement,morgan1994programming}. Fiat, Cozy, and
Primrose choose among the implementations a specification admits, Primrose
as annotations on Rust~\cite{delaware2015fiat,loncaric2016cozy,Qin2023Primrose}. STOKE
optimizes floating-point code against a stated error bound rather than
equivalence~\cite{schkufza2014fp}. This paper asks what that idea is, as a
compilation problem for ordinary systems code, and what an automated compiler
can actually do with it.

The answer has three layers.

\paragraph{The Ideal Problem.}
The developer states a functional contract \(\Phi_{\text{func}}\) over
observable behavior, including named implementation parameters such as a
capacity of at least 1000, together with a resource budget and a cost
objective. Programs that satisfy \(\Phi_{\text{func}}\) are legal, and
those that also satisfy the budget are feasible. The ideal compiler
returns a cheapest feasible legal program. That legal set contains every
as-if variant of a correct baseline, and more programs exactly where
\(\Phi_{\text{func}}\) leaves a decision open.

\paragraph{A Restricted Compiler.}
\textsc{IntentCompile} proposes candidates from a generator (rewrites and
parameter holes), admits those a sound verifier proves within a time budget,
and hands the winner to an unmodified production compiler. What it can use
is the verified fragment of the generated set, and it emits \(P_1\) or a
program cheaper on the supplied \(\mathcal{O}\).

\paragraph{An Instance.}
\textsc{IntC} is a C instance of \textsc{IntentCompile}. It takes a C
baseline, a contract written as parameter constraints plus CBMC harnesses,
and an optional benchmark.
It rewrites remainder, division, and multiplication by a parameter. It also
applies representation rewrites: pad-and-mask, which grows a private store
while keeping API capacity, and structure-of-arrays. Arithmetic rewrites drop
a sqrt from a disk test or a double accumulator from an inner product. It fills
parameter holes by search near the baseline, admits candidates a bounded model
checker proves, and selects by measurement. On the FIFO ring, that winner is
the 1024-cell mask, \(2.3\times\) on the named benchmark. What it claims of an admitted
program is \(\models_K\): the harness holds on every execution the bound
covers. It emits reviewable source so a developer can inspect an admitted
change as C, including a change that disagrees with the baseline on an
observable result.

\paragraph{Contributions.}
\begin{itemize}
  \item The compilation problem: contracts with named implementation
  parameters, four kinds of slack, and an ideal objective over the legal set
  (\S\ref{sec:ideal}).
  \item A restricted compiler whose search space is a generator, whose
  admission is a sound verifier, including the bounded judgment
  \(\models_K\), and whose backend is an unmodified production compiler
  (\S\ref{sec:compiler}).
  \item The \textsc{IntentCompile} pipeline and \textsc{IntC}, a C instance
  of it (\S\ref{sec:pipeline}, \S\ref{sec:implementation}).
  \item Evaluation of \textsc{IntC} (\S\ref{sec:evaluation}): whether it
  runs on extracts from real-world projects, whether as-if compilation and
  the internals behave as claimed, and whether ranking under the same
  \texttt{/O2} emits a cheaper program on the named driver or keeps
  \(P_1\). On the ring that driver is \(2.3\times\).
\end{itemize}

Section~\ref{sec:overview} walks through the ring buffer and what existing
tools can do with it. Sections~\ref{sec:ideal}--\ref{sec:compiler} are the two
theory layers. Sections~\ref{sec:pipeline}--\ref{sec:implementation} are the
algorithm and the artifact. Section~\ref{sec:evaluation} answers the three questions of \textsc{IntC}.
Section~\ref{sec:related} places the work. Section~\ref{sec:discussion}
names the risks that remain. Section~\ref{sec:conclusion} closes.

\section{Overview}
\label{sec:overview}

\subsection{Motivating Example}
\label{sec:motivating_example}

A network ingress path buffers packets in a ring so that bursts do not cause
drops. Bursts reach 800 packets, and the developer provisions 25\% headroom, so
the ring needs at least 1000 slots. The queue keeps two free-running 64-bit
counters. \texttt{tail} counts enqueued packets and \texttt{head} counts
dequeued ones. The queue holds \texttt{tail - head} packets. Each counter value
maps to a storage slot by taking it modulo the capacity.

\begin{lstlisting}[language=C]
#define CAPACITY 1000
uint64_t ring[CAPACITY];
uint64_t head = 0, tail = 0;

bool enqueue(const uint64_t pkt) {
    if (tail - head == CAPACITY) return false;
    ring[tail % CAPACITY] = pkt;
    ++tail;
    return true;
}
\end{lstlisting}

Every enqueue and dequeue computes a remainder by 1000. MSVC~19.29 at
\texttt{/O2} on x86-64 compiles that remainder to a reciprocal multiply:

\begin{lstlisting}[language={[x86masm]Assembler}]
; i % 1000
mov  rax, 442721857769029239
mul  rcx
mov  rax, rcx
sub  rax, rdx
shr  rax, 1
add  rax, rdx
shr  rax, 9
imul rax, rax, 1000
sub  rcx, rax
mov  rax, rcx
ret
\end{lstlisting}

A 1024-cell ring uses \texttt{i \% 1024}, which is the same as
\texttt{i \& 1023}, and the same compiler emits a mask:

\begin{lstlisting}[language={[x86masm]Assembler}]
; i % 1024
and  ecx, 1023
mov  eax, ecx
ret
\end{lstlisting}

Table~\ref{tab:motivating} splits the requirement, the source, what an as-if
compiler may emit, and what a contract admits. As-if will not rewrite 1000 to
1024, because the two queues disagree on the 1001st consecutive enqueue from
empty. In isolation, on an Intel Core i7-14700K, a dependent chain of slot
computations takes a median of 2.67\,ns per step with capacity 1000 and
0.39\,ns with capacity 1024, a factor of 6.8. When the computations are
independent, the factor is 3.1. These figures measure the slot computation
alone. Whole queues, where memory traffic dilutes the difference, are measured
in Section~\ref{sec:evaluation}, and the larger ring uses 24 more slots.

Both capacities meet a requirement of at least 1000 slots. Which is better
depends on the deployment. On a server where the path is latency-critical and
memory is plentiful, 1024 is the better choice. On a microcontroller whose
memory budget is exactly the 1000-slot footprint, the extra 24 slots do not
fit. The preference is specific to the target, so neither capacity is right
in the source as a universal constant.

\begin{table}
\caption{Columns are the requirement, the source \(P_1\), what MSVC
\texttt{/O2} emits under as-if, and a program the contract admits. As-if
must keep the remainder that is in the source. A contract with
\texttt{CAPACITY}~\(\ge 1000\) admits 1024.}
\label{tab:motivating}
\scriptsize
\begin{tabular}{p{0.22\linewidth}p{0.22\linewidth}p{0.22\linewidth}p{0.22\linewidth}}
\toprule
Requirement & Source \(P_1\) & MSVC \texttt{/O2} (as-if) & Contract admits \\
\midrule
\(\texttt{CAPACITY} \ge 1000\), FIFO, full at \texttt{CAPACITY} &
\texttt{\#define CAPACITY 1000}; \texttt{i \% CAPACITY} &
reciprocal multiply for \texttt{\% 1000}; 1001st enqueue fails &
\texttt{CAPACITY 1024}; \texttt{i \& 1023}; 1001st enqueue succeeds \\
full at 1000; backing store not in the API &
\texttt{CAPACITY} and \texttt{STORAGE} both 1000; \texttt{i \% STORAGE} &
keeps remainder-by-1000 &
\texttt{STORAGE 1024}, \texttt{CAPACITY 1000}; mask; enqueue still fails at 1000 \\
\bottomrule
\end{tabular}
\end{table}

That change is not the only way to use a power-of-two backing store. A library
can allocate 1024 slots and still report full at 1000, using the extra slots
only so that indexing is a mask. Folly's queue after September 2026 does
exactly that: it rounds storage up internally and keeps the caller-facing
capacity~\cite{folly_pcq}. The two designs use the same hardware trick and are different contracts: one
changes when enqueue fails, and the other does not.

\subsection{What Existing Tools Can Do}
\label{sec:existing}

\paragraph{Conventional Compilers and Translation Validation.}
The two queues that become full at different lengths are not observationally
equivalent. Translation validators such as Alive2~\cite{lopes2021alive2} check
refinement against the source, which allows extra freedom only where the source
has undefined behavior, and this queue has none on its intended inputs.
Alive2 rejects the 1024-cell ring against the 1000-cell source: the two
disagree on the 1001st enqueue from empty.

\paragraph{Compiler Flags.}
Production compilers already relax as-if under flags, and
\texttt{-ffast-math} is one example: it licenses reassociation and other
algebraic rewrites that IEEE-754 would forbid, and the waiver can be scoped
with pragmas. Those rewrites change observable numeric behavior, including
treating NaN and infinities as unused, sometimes with measurable
implications, but they do not let the compiler replace a 1000-cell remainder
with a 1024-cell mask.

\paragraph{Compiler Data-Layout Transformations.}
Production compilers do change data layout. IBM's XL compilers reshape arrays
of structures~\cite{zhao2007forma}, HP's compilers split and reorder
structures~\cite{hundt2006practical}, and automatic pool allocation regroups
heap objects~\cite{lattner2005pool}. These transformations must remain
invisible to the program. They are often abandoned when the analysis cannot
prove that~\cite{hundt2006practical}, and changing a logical capacity is
visible by construction, so it is outside their reach. Padding a private array while
preserving the API is closer to what they already do, and still requires
proving that the extra cells are never observed.

\paragraph{Annotations and Assumptions.}
C++'s \texttt{[[assume]]} states facts the compiler may rely on. A false
assumption is undefined behavior, and the C++26 contracts design leaves out an
assume semantic for that reason~\cite{p2899r1}. Recent Clang attributes permit
array-of-structures to structure-of-arrays conversions on annotated
code~\cite{radtke2025annotation}. Both are trusted rather than checked, and
neither states that a capacity may grow.

\paragraph{Scheduling Languages.}
Halide and Exo separate an algorithm from its schedule~\cite{ragan2013halide,
ikarashi2022exo}. The algorithm's meaning, including the size of the ring,
stays fixed.

\paragraph{Specification-Based Synthesis.}
Cozy synthesizes collection implementations from specifications, and Primrose
selects container implementations that satisfy stated
properties~\cite{loncaric2016cozy,Qin2023Primrose}. Fiat refines specifications
into verified implementations in Coq~\cite{delaware2015fiat}. These systems
accept freedom from a specification, in their own languages or libraries.
Primrose comes closest to existing code: it reads annotated Rust and selects a
library implementation. This paper starts from an existing C program, a
contract over its observable behavior, and that language's production compiler.

\subsection{Intent Compilation on the Example}
\label{sec:overview_intent}

The developer states the requirement instead of a single capacity. In
\textsc{IntC} that is a YAML file plus C harnesses. The contract has four pieces:
a constraint on named parameters, a functional harness, an optional resource
budget, and a measured objective. On this queue:

\begin{lstlisting}[language=C]
/* CAPACITY >= 1000 */
/* FIFO; enqueue fails iff the ring holds CAPACITY packets */
/* optional static-memory ceiling */
/* measured ns/op on a given target */
\end{lstlisting}

\textsc{IntentCompile} (\S\ref{sec:pipeline}) rewrites \texttt{e \% CAPACITY} to
\texttt{((e) \& (CAPACITY - 1))}, which requires a power-of-two capacity, searches
for the admissible power of two closest to 1000, and asks a bounded model
checker whether every sequence of at most \(K\) operations from an empty queue
satisfies the harness (the bound of \(\models_K\), \S\ref{sec:verifier}).
Under a server profile with no tight memory budget, 1024
is admitted and, if it measures faster, emitted as source. Under an 8016-byte
static-memory ceiling (the 1000-slot baseline's footprint), 1024 fails the
budget at 8208 bytes and the baseline is kept.
It also tries a pad-and-mask form: a \texttt{STORAGE} hole, a mask on that
store, and fullness still at \texttt{CAPACITY}. That candidate is the
unexposed-slack case of Section~\ref{sec:slack}. When the baseline already names
both parameters, the mask rewrite on \texttt{STORAGE} is enough.

\section{The Ideal Problem}
\label{sec:ideal}

This section defines the compilation problem: a contract over observable
behavior, the programs that satisfy it, and an objective that picks a cheapest
one. Leroy's preservation of a
specification~\cite[\S2.1]{leroy2009backend} and the refinement
calculi~\cite{back1998refinement,morgan1994programming} are the setting. Named
implementation parameters, a taxonomy of slack, that objective, and then a
split between the functional contract and the optimization contract make the
problem specific enough to compile against.

The definitions apply to any language with a deterministic sequential semantics
and a notion of observable behavior. The running examples are in C, and the
as-if rule and undefined behavior are the same in C++.

\subsection{Programs and the As-If Rule}
\label{sec:setting}

Let \(\Sigma\) be the set of initial states of a program, covering its inputs
and initial memory. Running program \(P\) on \(\sigma \in \Sigma\) has one of
three results. It may terminate in a final state \(\sigma'\) after producing a
trace \(\tau \in \mathcal{T}\) of observable events, written
\(\langle P, \sigma \rangle \Downarrow (\sigma', \tau)\). It may run forever,
written \(\mathsf{div}\). Or it may reach an operation whose meaning the
language leaves undefined, written \(\mathsf{ub}\). We write
\(\llbracket P \rrbracket(\sigma)\) for the result.

In C and C++, the observable behavior consists of I/O and accesses to volatile
objects~\cite[\S4.1.2]{n4950}. Local values, the order of intermediate
computations, and unescaped layout are not observable, so comparisons use only
those observable parts.

\begin{definition}[Observational Equivalence]
\(P_1 \approx P_2\) iff \(\llbracket P_1 \rrbracket(\sigma)\) and
\(\llbracket P_2 \rrbracket(\sigma)\) have the same observable part for every
\(\sigma \in \Sigma\). We write \([P]_\approx\) for the set of programs
observationally equivalent to \(P\).
\end{definition}

\([P]_\approx\) is the space a conventional optimizing compiler may search.
Undefined behavior loosens this slightly: on inputs where the source is
undefined, the compiled program may do anything~\cite{lopes2021alive2,
lee2017taming}.

\subsection{Named Parameters and Functional Contracts}
\label{sec:contract}

Implementation parameters are named constants in the source, written
\(\theta_P\). The set of names is part of the contract. Extracting
\(\theta_P\) is syntactic: a \texttt{\#define} capacity, not a property of
\(\llbracket P \rrbracket\). Two observationally equivalent programs may
differ in \(\theta\). The contract is indexed by \(\theta\), so
\(\mathit{Post}_\theta\) can depend on the capacity the implementation chose.

\begin{definition}[Functional Intent Contract]
\label{def:contract}
A functional intent contract is a triple
\[
\Phi_{\text{func}} = \langle \mathit{Pre}, \mathit{Post}, \mathcal{I} \rangle
\]
where \(\mathit{Pre} \subseteq \Sigma\) is the set of initial states the
program must handle, \(\mathcal{I}\) is the set of acceptable parameter
valuations, and \(\mathit{Post}\) gives, for each \(\theta \in \mathcal{I}\),
a relation \(\mathit{Post}_\theta\) between an initial state and the
acceptable observable results.
Program \(P\) satisfies \(\Phi_{\text{func}}\), written
\(P \models \Phi_{\text{func}}\), iff \(\theta_P \in \mathcal{I}\) and
\[
\forall \sigma \in \mathit{Pre}.\ \exists \sigma', \tau.\
\langle P, \sigma \rangle \Downarrow (\sigma', \tau)
\ \wedge\ (\sigma, (\sigma', \tau)) \in \mathit{Post}_{\theta_P}.
\]
\end{definition}

The triple follows Design by Contract~\cite{meyer1988eiffel}.
Satisfaction requires termination without undefined behavior on every input
in \(\mathit{Pre}\), and the contract mentions only observable results. The
set of satisfying programs is
\(\mathrm{Sat}(\Phi_{\text{func}}) = \{P \mid P \models \Phi_{\text{func}}\}\).
We write \(P_1 \sim_{\Phi_{\text{func}}} P_2\) when both lie in that set.

Each component states one kind of intent. \(\mathit{Pre}\) states which
inputs matter. \(\mathcal{I}\) states acceptable parameter valuations, so a
queue that must absorb bursts of 800 packets with 25\% headroom has
\(\mathcal{I} = \{\theta \mid \theta_{\mathrm{cap}} \ge 1000\}\).
\(\mathit{Post}_\theta\) states what the program must do, relative to the
parameters it chose: enqueued packets come out in order, and an enqueue fails
exactly when \(\theta_{\mathrm{cap}}\) packets are stored.

\begin{definition}[Realizability]
\(\Phi_{\text{func}}\) is vacuous if \(\mathit{Pre} = \emptyset\). It is
contradictory if no parameter valuation, input, and result jointly satisfy
\(\mathit{Post}\). It is realizable if there exists \(\theta \in \mathcal{I}\)
such that every \(\sigma \in \mathit{Pre}\) has at least one acceptable
result under \(\mathit{Post}_\theta\).
\end{definition}

Vacuous contracts are rejected before search: they have empty
\(\mathit{Pre}\). A contract whose \(\mathit{Post}\) no program can satisfy
is discovered later, when no candidate is proved. Realizability is the
\(\forall\sigma\,\exists o\) condition.

\subsection{Slack}
\label{sec:slack}

\(\mathrm{Sat}(\Phi_{\text{func}})\) differs from \([P]_\approx\) where the
contract leaves a decision open. There are four ways to leave one open.

\begin{itemize}
  \item \emph{Input slack.} \(\mathit{Pre} \neq \Sigma\). On excluded inputs
  any behavior is acceptable. This is the kind undefined behavior already
  provides, and the kind SeGaBench grants by fixing a valid-input
  domain~\cite{jiang2026segabench}.
  \item \emph{Output slack.} For some covered input, \(\mathit{Post}_\theta\)
  accepts two distinct observable results. A process-ID allocator must return
  an unused ID, or \(\texttt{NONE}\) when every ID is in use. Scanning from
  ID 1, and continuing from a cursor, both satisfy that, and they return
  different IDs once an ID below the cursor has been freed. No conventional
  compiler may replace one by the other.
  \item \emph{Logical-parameter slack.} The contract accepts several
  parameter valuations, and the choice is visible in the results. Queues of
  capacity 1000 and 1024 reject a full-queue enqueue at different points. Both
  are acceptable when \(\mathcal{I}\) requires only a capacity of at least
  1000.
  \item \emph{Unexposed slack.} Extra resources or a different representation
  that the interface does not expose. A queue that allocates 1024 slots and
  still becomes full at 1000 uses a power-of-two mask while preserving
  caller-facing capacity. Folly and DPDK's exact-size rings are this
  case~\cite{folly_pcq,dpdk_rte_ring}. If the extra cells and the mask are
  unobservable at the API, the as-if rule may already permit the change,
  provided the analysis can see it. The contract still has to license the
  extra memory, which is a budget concern
  (\S\ref{sec:opt_contract}).
\end{itemize}

Input and output slack are don't-cares in Karnaugh's sense: an incompletely
specified function can be minimized by assigning combinations the
specification leaves open~\cite{karnaugh1953map}.
Logical-parameter slack and unexposed slack concern the implementation
rather than unspecified I/O. The first is a named parameter visible in the
results, and the second is extra store or a representation the interface
does not name.

The four kinds of slack are properties of \(\Phi_{\text{func}}\). Whether
\(\mathrm{Sat}(\Phi_{\text{func}})\) is larger than \([P]_\approx\) is
settled by local patching.

A language supports \emph{local patching} if for every program \(P\), state
\(\sigma_0\), and producible observable result \(o\), some program behaves
like \(P\) everywhere except at \(\sigma_0\), where it produces \(o\). Any
language that can test its input against a constant has this property,
including essentially every general-purpose language in use today.

\begin{proposition}[Legal Set Expands Exactly on Observable Slack]
\label{prop:expansion}
Assume local patching, and that each \(\mathit{Post}_\theta\) with
\(\theta \in \mathcal{I}\) is realized by some program. Let
\(P \models \Phi_{\text{func}}\). Then
\(\mathrm{Sat}(\Phi_{\text{func}})\) contains a program that is not
observationally equivalent to \(P\) iff the contract has input, output, or
logical-parameter slack, counting only results some program produces.
Unexposed slack is not in that list: it need not change those results.
\end{proposition}

\noindent\emph{Proof sketch.} If the contract accepts a result on some input
that differs from \(P\)'s, patching produces a legal program outside
\([P]_\approx\). Without slack, every covered input has one acceptable result,
the same for every allowed parameter, and every input is covered, so every
legal program matches \(P\). (Appendix~\ref{app:proofs}.)

When the contract has no input, output, or logical-parameter slack, intent
compilation is conventional compilation: the approach generalizes the as-if
rule rather than replacing it. Unexposed slack can still change cost and
memory. The proposition does not say that any extra legal program is faster.

\subsection{The Ideal Objective}
\label{sec:ideal_obj}

Let \(\mathcal{L}_{\text{Target}}\) be an unmodified production compiler for
the target. It implements the as-if rule: idealized, it searches
\([P]_\approx\). Let \(\mathcal{O}(\mathrm{Bin}, \text{Target})\) rank
binaries by cost on that target. A conventional compiler, idealized,
returns a cheapest binary in the as-if class of \(P_1\):
\begin{equation}
\label{eq:traditional}
\mathrm{Bin}^*_{\text{trad}} \in
\arg\min_{\mathrm{Bin} \in \{\mathcal{L}_{\text{Target}}(P') \mid
P' \in [P_1]_\approx\}}
\mathcal{O}(\mathrm{Bin}, \text{Target}).
\end{equation}
An intent compiler wants the same ranking over every binary that
satisfies the contract:
\begin{equation}
\label{eq:intent_pure}
\mathrm{Bin}^*_{\text{intent}} \in
\arg\min_{\mathrm{Bin} \models \Phi_{\text{func}}}
\mathcal{O}(\mathrm{Bin}, \text{Target}).
\end{equation}
No compiler searches either set as a whole. Conventional compilation
applies heuristic rewrites inside \([P_1]_\approx\), the passes of
\(\mathcal{L}_{\text{Target}}\), and may never reach
\(\mathrm{Bin}^*_{\text{trad}}\). Intent compilation has the same gap on a
larger set, and splits~\eqref{eq:intent_pure} into a source step and an
as-if step:
\begin{align}
\label{eq:intent_source}
P^\star &\in \arg\min_{P' \models \Phi_{\text{func}}}
\mathcal{O}(\mathcal{L}_{\text{Target}}(P'), \text{Target}), \\
\label{eq:intent_lower}
\mathrm{Bin}^*_{\text{intent}} &\in
\arg\min_{\mathrm{Bin} \in \{\mathcal{L}_{\text{Target}}(P'') \mid
P'' \in [P^\star]_\approx\}}
\mathcal{O}(\mathrm{Bin}, \text{Target}).
\end{align}
The first level rewrites the source toward a cheaper program in
\(\mathrm{Sat}(\Phi_{\text{func}})\), and the second is~\eqref{eq:traditional}
with \(P^\star\) in place of \(P_1\). That composition does not guarantee a
minimum of~\eqref{eq:intent_pure}, but when \(P_1\) satisfies the contract it
is never worse than as-if compilation of \(P_1\), because the first level
searches a set that contains every as-if variant of \(P_1\). It separates
program-to-program editing under the contract from semantic lowering of a
single source, and the second level still yields a contract-satisfying
binary from a contract-satisfying source.

\begin{proposition}[As-If Programs Remain Legal]
\label{prop:refine}
If \(P \models \Phi_{\text{func}}\), \(P' \approx P\), and
\(\theta_{P'} = \theta_P\), then \(P' \models \Phi_{\text{func}}\).
\end{proposition}

\noindent\emph{Proof sketch.} A violation by \(P'\) would be a bad parameter
or a bad observable result on some admissible input. The parameters match, so
the first case is impossible. Observational equivalence transfers the second
case to \(P\), contradicting \(P \models \Phi_{\text{func}}\).
(Appendix~\ref{app:proofs}.)

The side condition \(\theta_{P'} = \theta_P\) is required because
\(\mathit{Post}\) is indexed by \(\theta\). Hence the second level
preserves the contract: if \(P^\star \models \Phi_{\text{func}}\), every
program in \([P^\star]_\approx\) with the same \(\theta\) does too.

\subsection{Optimization Contracts}
\label{sec:opt_contract}

The source step~\eqref{eq:intent_source} and the as-if
step~\eqref{eq:intent_lower} are not ranked by the same objective, even
though both displays used \(\mathcal{O}\). The as-if step is
\(\mathcal{L}_{\text{Target}}\)'s search inside \([P^\star]_\approx\).
The source step compares programs that need not be observationally
equivalent, which that backend does not rank. A ranking
for~\eqref{eq:intent_source} has to be supplied as part of the intent.

That ranking searches a larger set than \([P_1]_\approx\). A rewrite in
that set can use more resources than \(P_1\), exceeding what the developer
allocated for the program, for instance a larger memory footprint. Those
allocations are implicit in \(P_1\), so as-if cannot raise them and
\(\mathrm{Sat}(\Phi_{\text{func}})\) does not inherit them. Therefore, an
operational budget \(\mathcal{B}\) is also needed, stating which
implementations are valid to search: those whose resource use stays within
what the developer is willing to allocate to the program.

\begin{definition}[Legality and Feasibility]
\label{def:legal}
A program \(P\) is \emph{legal} iff \(P \models \Phi_{\text{func}}\), and
\emph{feasible} iff \(P \models \mathcal{B}\).
\end{definition}

Search is over programs that are legal and feasible.
\(\mathcal{B}\) and \(\mathcal{O}\) stay outside \(\Phi_{\text{func}}\)
because legality, feasibility, and ranking are different jobs.
\(\Phi_{\text{func}}\) states which programs are legal, \(\mathcal{B}\)
which of those are feasible, and \(\mathcal{O}\) ranks them. That ranking
is supplied as part of the intent. It may be computed from the program
text or measured on a profiled run. \textsc{IntC} uses wall-clock time on
a named driver. Cyclomatic complexity occupies the same slot
(\S\ref{sec:related}). A missed latency target is then suboptimal rather
than illegal, and folding the target into the functional contract would
reject a legal baseline. Budgets describe the machine and change with the
deployment, while \(\Phi_{\text{func}}\) does not.
One functional contract can then yield a 1024-slot ring on a server and a
1000-slot ring under a tight memory budget. A check of \(\mathcal{B}\) on source
approximates a constraint on the compiled binary: padding, alignment, and
allocator overhead are decided later.

Adding \(\mathcal{B}\) as a constraint on the source
step~\eqref{eq:intent_source} gives
\begin{equation}
\label{eq:intent}
P^\star \in \arg\min_{\substack{P' \models \Phi_{\text{func}} \\ P' \models \mathcal{B}}}
\mathcal{O}(\mathcal{L}_{\text{Target}}(P'), \text{Target}),
\qquad
\mathrm{Bin}^*_{\text{intent}} = \mathcal{L}_{\text{Target}}(P^\star),
\end{equation}
where the second equality is~\eqref{eq:intent_lower} once
\(\mathcal{L}_{\text{Target}}\) has searched \([P^\star]_\approx\), and
Section~\ref{sec:compiler} searches a generated subset of that source set.

\begin{definition}[Intent Specification]
\label{def:intent_spec}
An intent specification is a triple
\[
\Phi = \langle \Phi_{\text{func}}, \mathcal{B}, \mathcal{O} \rangle
\]
where \(\Phi_{\text{func}}\) is a functional intent contract
(Definition~\ref{def:contract}), \(\mathcal{B}\) is a conjunction of
operational constraints on a candidate (static memory, a ban on dynamic
allocation, alignment), and \(\mathcal{O}\) is a cost ranking of binaries on
a target.
\end{definition}

If \(P_1\) already satisfies the contract and the budget, the minimum
in~\eqref{eq:intent} is never worse than the one in~\eqref{eq:traditional}
among programs that keep \(\theta_{P_1}\), by Proposition~\ref{prop:refine}.

\section{A Restricted Intent Compiler}
\label{sec:compiler}

An automated compiler cannot evaluate~\eqref{eq:intent}. Membership in
\(\mathrm{Sat}(\Phi_{\text{func}})\) is undecidable for interesting
contracts, by Rice's theorem~\cite{rice1953classes}, and even a decidable
proof still has a time budget. This section replaces the source
step~\eqref{eq:intent_source} with a generator, a sound verifier \(V\) with
budget \(t\), and a fallback. The as-if step~\eqref{eq:intent_lower} remains
\(\mathcal{L}_{\text{Target}}\). The contract bounds what the compiler may
do, and what \(V\) proves within \(t\) bounds what it can use.

\subsection{Generator}
\label{sec:generator}

Let \(R\) be a catalog of source-to-source rewrites, each with a syntactic
applicability condition and optional side constraints on parameters, for
example that capacity is a power of two. Let \(D\) be a finite domain for
parameter holes, as machine integers are. \(R\) and \(D\) are inputs of
\(\mathrm{Gen}\). The \emph{generated set} is
\[
\mathrm{Gen}(P_1, R, D)
= \{ S[h] \mid S \text{ is obtained from \(P_1\) by a finite sequence of
rewrites from \(R\)},\ h \in D \}.
\]
A program is in \(\mathrm{Gen}\) only if a finite sequence of rewrites from
\(R\) produces its structure and some \(h \in D\) fills its holes, so a
cursor-based PID allocator, a linked queue, or a ring that pads storage
while keeping caller-facing capacity appears only if \(R\) contains that
rewrite. Rewrites may change observable behavior, and their outputs are
candidates until verified.

\subsection{Verifier}
\label{sec:verifier}

A verifier \(V\) takes a program, a contract, and a time budget \(t\) and
answers \(\mathsf{proved}\), a counterexample, or \(\mathsf{unknown}\). It
is \emph{sound} if \(\mathsf{proved}\) implies
\(P \models \Phi_{\text{func}}\) under an exact model of machine arithmetic,
including overflow.

Many useful verifiers are bounded. A bound \(K\) (operations, unwindings)
induces a judgment \(P \models_K \Phi_{\text{func}}\): the contract holds on
every execution the harness covers up to \(K\), and a sound bounded verifier
may conclude \(\mathsf{proved}\) only for that judgment. Timeouts, solver
failures, exhausted unwindings, and a property failure on a truncated path
are \(\mathsf{unknown}\), never \(\mathsf{proved}\).

There are two ways to certify a candidate \(P_2\). A relational check proves
that \(P_2\) behaves like \(P_1\), taking observational equivalence to
\(P_1\) as the bound. The
developer writes no \(\Phi_{\text{func}}\) and the query contains both
programs. Every candidate outside \([P_1]_\approx\) is then rejected.
That query inherits
\(P_1\)'s complexity and may be undecidable, but it is cheap when the
programs are structurally close. Those checks remain the right tool for a
small patch when no contract is available~\cite{lopes2021alive2}. A
decoupled check proves \(P_2 \models \Phi_{\text{func}}\) (or
\(\models_K\)) directly, taking \(\Phi_{\text{func}}\) as the
specification. The developer must write that contract, and the query
contains only the candidate. This check does not inherit \(P_1\)'s complexity, but the
candidate itself has to lie in a fragment \(V\) can decide, or decide up to
\(K\). \(P_1\) need not satisfy the contract: it is a starting point for the
generator and a fallback. Search stays near that baseline so a candidate is
still a program the developer can read against the sources they already
have.

This compiler uses the decoupled check, so the \emph{admitted set} at budget
\(t\) is
\[
\mathcal{A}(P_1, \Phi, t)
= \{ P_2 \in \mathrm{Gen}(P_1, R, D) \mid
V(P_2, \Phi_{\text{func}}, t) = \mathsf{proved}
\ \text{and}\ P_2 \text{ passes the check for }\mathcal{B} \}.
\]
The two conjuncts are different checks: \(V\) admits on \(\Phi_{\text{func}}\),
and a provided checker admits on \(\mathcal{B}\), not a second run of \(V\).
When \(V\) is unbounded and sound, \(\mathcal{A}\) is a subset of
\(\mathrm{Sat}(\Phi_{\text{func}})\). When \(V\) is bounded, it is a subset
of the \(K\)-satisfying programs. Only programs in \(\mathrm{Gen}\) are
submitted to \(V\), and \(\mathsf{unknown}\) shrinks \(\mathcal{A}\) without
unsounding an emit. A
bounded model checker or another sound oracle can play \(V\), and they
differ in which contracts they express and how much of \(\mathrm{Gen}\) they
reach in \(t\). A synthesizer built on \(V\) can produce only what \(V\) can
prove~\cite{Srivastava2010VS3}.

\subsection{Backend}
\label{sec:backend}

The as-if step~\eqref{eq:intent_lower} is unchanged: the restricted compiler
hands a source program to \(\mathcal{L}_{\text{Target}}\).

\paragraph{Compiler Correctness Assumption.}
For every program \(P\) and every \(\sigma\) on which \(P\) has no undefined
behavior, \(\llbracket \mathcal{L}_{\text{Target}}(P) \rrbracket(\sigma)\)
and \(\llbracket P \rrbracket(\sigma)\) have the same observable part. That
is the as-if rule as compilers implement it, trusted rather than verified.

Proposition~\ref{prop:refine} already says that an as-if program with the
same \(\theta\) remains legal. \(\mathcal{L}_{\text{Target}}\) is that as-if
on inputs where the source has no undefined behavior, and it keeps
\(\theta\), so the unbounded case of the next lemma is that proposition
instantiated at \(P' = \mathcal{L}_{\text{Target}}(P^\star)\). The bounded
judgment is the same fact on the traces \(\models_K\) covers.

\begin{lemma}[Lowering Preserves a Proved Contract]
\label{lem:lowering}
If \(P^* \models \Phi_{\text{func}}\) (respectively \(\models_K\)), then
\(\mathcal{L}_{\text{Target}}(P^*) \models \Phi_{\text{func}}\)
(respectively \(\models_K\)).
\end{lemma}

\noindent\emph{Proof sketch.} Implementation parameters remain those of the
source. \(P^*\) has no undefined behavior on inputs in \(\mathit{Pre}\) (or
on the traces \(\models_K\) covers), so the compiler correctness assumption
applies, and Proposition~\ref{prop:refine} gives the unbounded case.
(Appendix~\ref{app:proofs}.)

The lemma needs a contract over observable results and a candidate that is
UB-free on the covered inputs, so the trusted base is the production
compiler, the verifier, and the encoding of the language into the verifier's
logic, which is larger than that of a verified compiler or of Fiat's
end-to-end proofs~\cite{leroy2009backend,delaware2015fiat}.

\subsection{Selection and Fallback}
\label{sec:selection}

Winner selection lowers the baseline and every admitted candidate and keeps
the binary with the lowest \emph{measured} cost:
\begin{equation}
\label{eq:winner}
\mathrm{Bin}^* = \arg\min_{\mathrm{Bin} \in \{\mathcal{L}_{\text{Target}}(P)
\mid P \in \mathcal{A} \cup \{P_1\}\}}
\mathcal{O}(\mathrm{Bin}, \text{Target}),
\end{equation}
preferring \(P_1\) on ties. Each admitted source, and \(P_1\), is compiled
with \(\mathcal{L}_{\text{Target}}\) and ranked by measured cost on the
profiled workload and hardware, and the emitted artifact is the winning
source, not an object file.

If \(\mathcal{A}\) is empty, the compiler emits \(P_1\) unchanged. A change
is admitted only after a proof of \(\Phi_{\text{func}}\) (or \(\models_K\)).
Callers whose tests depended on incidental behavior the contract left open
can still observe a difference.

\subsection{What the Restricted Compiler Guarantees}
\label{sec:restricted_thm}

\begin{theorem}[Soundness of a Restricted Intent Compiler]
\label{thm:restricted}
Let \(V\) be sound for \(\models\) or for \(\models_K\), and let the compiler
emit source \(P^*\) and binary
\(\mathrm{Bin}^* = \mathcal{L}_{\text{Target}}(P^*)\).
\begin{enumerate}
  \item If \(P^* \neq P_1\), then \(P^*\) is in \(\mathrm{Gen}(P_1, R, D)\),
  \(V\) proved it, and it passed \(\mathcal{B}\). Hence
  \(P^* \models \Phi_{\text{func}}\) or \(P^* \models_K \Phi_{\text{func}}\),
  matching \(V\), and the same judgment holds of \(\mathrm{Bin}^*\) by
  Lemma~\ref{lem:lowering}.
  \item The measured cost of \(\mathrm{Bin}^*\) is at most that of
  \(\mathcal{L}_{\text{Target}}(P_1)\), when measurement runs.
  \item If \(P_1 \models \Phi_{\text{func}}\) (or \(\models_K\)), then
  \(\mathrm{Bin}^*\) satisfies the same judgment in every case.
\end{enumerate}
\end{theorem}

\noindent\emph{Proof sketch.} Winner selection returns \(P_1\) or a member
of \(\mathcal{A}\). Membership in \(\mathcal{A}\) is part 1. Part 2 is
minimization over a set that contains the compiled baseline. Part 3 uses
part 1 or Lemma~\ref{lem:lowering} on \(P_1\).
(Appendix~\ref{app:proofs}.)

If \(P_1\) already satisfies the contract, every emit does. If it does not,
the compiler replaces it only with a proved program or leaves it as written,
so the emit is never less safe than compiling \(P_1\) with
\(\mathcal{L}_{\text{Target}}\) alone. If some generated hole assignment of
a given sketch satisfies the contract and \(V\) decides every query, a
baseline-directed search that enumerates \(D\) in order of distance to
\(\theta_{P_1}\) returns a closest satisfying assignment
(Theorem~\ref{thm:cegis}). If no generated program satisfies the contract,
the compiler keeps \(P_1\), and a different structure requires a rewrite in
\(R\).

The compiler ranks by \(\mathcal{O}\) among programs \(V\) has proved.
How \(\mathrm{Gen}\) is explored, including any score that orders rewrites,
does not change that ranking. Verification cost and execution cost need not
agree~\cite{wagner2013overify,rinard2026quantitative}, so the fastest program
may lie outside \(\mathcal{A}\). They agree on the ring: integer division is
expensive on the processor and expensive to
bit-blast~\cite{niemetz2024bitblasting}, and the rewrite that removes the
remainder also simplifies the proof.

\section{The \textsc{IntentCompile} Pipeline}
\label{sec:pipeline}

\textsc{IntentCompile} takes a baseline \(P_1\), a specification
\(\Phi = \langle \Phi_{\text{func}}, \mathcal{B}, \mathcal{O} \rangle\), and a target. It
builds \(\mathcal{A}\) and then applies~\eqref{eq:winner}.
Appendix~\ref{app:algorithms} gives each stage as an algorithm.

\begin{enumerate}
  \item \textbf{\textsc{CheckRealiz}} (Algorithm~\ref{alg:checkrealiz}) checks that \(\mathcal{I}\) is
  satisfiable, and reports an unsatisfiable core on conflict. The stage is
  one solver call and returns when that call does. A timeout is
  \(\mathsf{unknown}\) and stops the pipeline. A harness that no program can
  satisfy is discovered later, when no candidate is proved.

  \item \textbf{\textsc{TheoryLower}} (Algorithm~\ref{alg:theorylower})
  applies rewrites from \(R\) that strictly decrease a nonnegative integer
  score on program structure. The score is a weighted count of constructs
  that make the query expensive for \(V\), so descent aims at a sketch the
  checker can decide, and the loop terminates because each rewrite drops
the score. Crossing the threshold stops that descent. A sketch is still
handed to CEGIS when a rewrite applied, even if the score remains above
the threshold. If no rewrite applied, there is no TheoryLower seed and
\textsc{GeneratePool} still starts from \(P_1\). Outputs are sketches with holes and side
constraints. \textsc{IntC}'s weights are in
Section~\ref{sec:implementation}.

  \item \textbf{\textsc{Baseline-Directed CEGIS}} (Algorithm~\ref{alg:cegis}) fills holes. An optimizer
  proposes valuations in \(\mathcal{I}\) that meet the side constraints and
  are as close as possible to \(\theta_{P_1}\). Holes not in the baseline
  are minimized subject to those constraints. \(V\) checks each candidate.
  A counterexample refutes that valuation, failing inputs are not constraints
  on later proposals, and an \(\mathsf{unknown}\) result stops the search. A
  finite attempt cap preserves soundness and may miss a farther satisfying
  assignment. The loop is enumerative CEGIS on a finite
  hole domain. The same loop fills TheoryLower sketches and GeneratePool
sketches. The sketch does not come from the developer, and hole-filling
distance is to the baseline so filled holes stay near the program the
developer already
wrote~\cite{solarlezama2006sketch, dantoni2016qlose}.

  \item \textbf{\textsc{BudgetNormalize}} (Algorithm~\ref{alg:budgetnormalize}) does not generate programs. It
  applies the provided checker for \(\mathcal{B}\) to the CEGIS seeds
  and drops those that miss \(\mathcal{B}\). If none remain,
  \textsc{GeneratePool} still runs from \(P_1\). \(P_1\) is the fallback of
  winner selection and is not re-verified.

  \item \textbf{\textsc{GeneratePool}} (Algorithm~\ref{alg:generatepool}) adds further proved programs. It
  starts from \(P_1\), not from TheoryLower's seeds, so it still runs when
  TheoryLower is stuck or CEGIS returns nothing. It applies a second catalog
  of representation rewrites, fills them by CEGIS, and keeps those that
  \(V\) proves and that fit \(\mathcal{B}\).

  \item \textbf{Winner selection} (Algorithm~\ref{alg:winner}) is the ranking by~\eqref{eq:winner}
  when the contract names a benchmark. Ranking runs only on \(P_1\) and on
  programs \(V\) proved.
  When there is no benchmark, the first admitted candidate is emitted
  (Section~\ref{sec:implementation}). The stage emits the chosen program as
  source, with the proof bound and the measurements, because an admitted
  change can disagree with \(P_1\) on an observable result and that
  disagreement has to be readable as source.
\end{enumerate}

\begin{theorem}[Properties of Baseline-Directed CEGIS]
\label{thm:cegis}
Let \(S[h]\) range over a finite hole domain \(D\), and let \(N\) be the
attempt cap (\(N = |D|\) if there is no cap).
\begin{enumerate}
  \item If CEGIS returns \(S[h]\), then \(V\) proved \(S[h]\).
  \item If \(V\) returns on each query, CEGIS returns after at most
  \(\min(N, |D|)\) iterations.
  \item If some \(h^\circ \in D\) makes \(S[h^\circ]\) satisfy the contract
  (or \(\models_K\)), \(V\) never answers \(\mathsf{unknown}\), and there is
  no attempt cap, then CEGIS returns a satisfying assignment closest to
  \(\theta_{P_1}\).
\end{enumerate}
\end{theorem}

\noindent\emph{Proof sketch.}
Return is only on \(\mathsf{proved}\). \(V\) respects \(t\), each miss
removes one element of \(D\), and the cap is finite, so the loop returns.
A satisfying assignment is never excluded, so the first proved valuation is
a closest one, unless an attempt cap stops first.
(Appendix~\ref{app:proofs}.)

The bound \(|D|\) is a worst case. Side constraints such as ``is a power of
two'' shrink it. An attempt cap preserves soundness and may drop part 3.

On the ring buffer, \(\mathcal{I}\) requires capacity at least 1000.
\textsc{TheoryLower} produces a mask sketch with a power-of-two side
constraint. CEGIS's closest value is 1024. The verifier sees the masked
program and the harness. The remainder in \(P_1\) does not appear in the
query.

\section{Implementation}
\label{sec:implementation}

\textsc{IntC} is a Python source-to-source compiler for C. Tree-sitter
locates spans in the C text, and there is no intermediate representation.
Z3 solves parameter constraints, CBMC 6.11.0 is \(V\), and MSVC 19.29
compiles and measures. The cases and the scripts that produce the numbers
in Section~\ref{sec:evaluation} are available with the paper.

\paragraph{Contracts.}
Implementation parameters are \texttt{\#define} constants.
\(\mathcal{I}\) lists constraints over those names, written
\texttt{CAPACITY >= 1000}.
\(\mathit{Pre}\) and \(\mathit{Post}\) are harnesses in C. A harness calls
the program with arbitrary inputs, restricts them with assumptions, and
asserts observable results~\cite{clarke2004cbmc}. A YAML file records
\(\mathcal{I}\), \(\mathcal{B}\), \(\mathcal{O}\), and the paths of those
harnesses, and does not generate the asserts. Unwind bounds and timeouts
configure \(V\) and are not \(\mathit{Post}\). The harness sees the
candidate's \(\theta\) through a generated header, which is how
\(\mathit{Post}_\theta\) depends on \(\theta\). An optional
compiled-\texttt{sizeof} ceiling and a named benchmark complete
\(\mathcal{B}\) and \(\mathcal{O}\).

\paragraph{Verifier.}
CBMC checks each harness together with out-of-bounds accesses, invalid
pointers, signed overflow, division by zero, and invalid shifts. Loops are
unrolled with unwinding assertions. A bound that is too small yields
\(\mathsf{unknown}\), not a proof, including when a later property also
fails on that truncated path. Arrays use the theory of arrays
(\texttt{--arrays-uf-always}). Without that encoding, CBMC expands a
1024-element ring element by element and does not finish within 120 seconds
even for three operations. Admission is \(\models_K\) for the harness bound
\(K\).

\paragraph{Stages.}
The stages are those of Section~\ref{sec:pipeline}. Only
\textsc{CheckRealiz} stops the pipeline: a failed TheoryLower, an empty
CEGIS, or an empty budget filter still runs \textsc{GeneratePool} from
\(P_1\) and then winner selection.
\begin{enumerate}
  \item \textbf{\textsc{CheckRealiz}} checks \(\mathcal{I}\) with Z3.
  Unknown on \(\mathcal{I}\) fails the stage, as does a name in
  \(\mathcal{I}\) that is not a \texttt{\#define}.

  \item \textbf{\textsc{TheoryLower}}'s score is
  \(10\cdot n_{\mathrm{rem/div}} + 3\cdot n_{\mathrm{mul}} + 2\cdot n_{\mathrm{loop}} + n_{\mathrm{deref}}\),
  where the four counts are remainders and divisions that are not powers
  of two, variable--variable multiplies, loops, and pointer dereferences.
  Power-of-two remainders and divisions, and multiplies by literals, are
  omitted. The sum is a nonnegative integer, so a rewrite that drops at
  least one counted construct terminates. The contract supplies the stop
threshold, default 2. Crossing it stops that descent. The catalog is three
rewrites, and those
  rewrites introduce the holes.
  \begin{itemize}
    \item \texttt{e \% N} becomes \texttt{((e) \& (N - 1))} when \(N\) is a
    parameter, with side constraint \(\mathrm{is\_pow2}(N)\).
    \item \texttt{e / D} becomes \texttt{((e) >> S)} when \(D\) and \(S\) are
    parameters, with \(\mathrm{pow2\_pair}(D,S)\).
    \item \texttt{e * M} becomes a parenthesized left shift of \(e\) by a
    fresh hole \(S_M\) with \(\mathrm{pow2\_pair}(M,S_M)\).
  \end{itemize}
  CEGIS sees a TheoryLower sketch only when one of those applied.

  \item \textbf{Baseline-directed CEGIS} asks Z3's optimizer for the
  closest unrefuted valuation of baseline parameters, minimizes extra
  holes subject to the side constraints, then asks CBMC. A counterexample
  excludes that valuation and is not a constraint on later proposals. A
  timeout or an undecided harness is \(\mathsf{unknown}\): CEGIS stops
  and does not add that valuation to the excluded set. The attempt cap is
  a contract bound whose default is 16.

  \item \textbf{\textsc{BudgetNormalize}} sums compiled \texttt{sizeof} of
  file-scope objects, drops a candidate that exceeds the ceiling, and
  treats a missing ceiling as no check.

  \item \textbf{\textsc{GeneratePool}} starts from \(P_1\). Each matching
  rewrite becomes one sketch \(r(P_1)\), filled by the same CEGIS loop.
  The catalog is general representation rewrites. It is extended with one
  domain formula, registered the same way a developer would add a formula
  from their domain.
  \begin{itemize}
    \item Representation:
    \begin{itemize}
      \item pad-and-mask (a larger backing array so remainder indexing
      becomes a mask)
      \item narrowing a 32-bit space tag to \texttt{uint8\_t}
      \item array-of-structs to structure-of-arrays
      \item aligning a loop bound to a multiple of 4
      \item predicating a sparse \texttt{if (flag[i])} update
      \item 1D halo pad (ghost cells so a 3-point stencil loop does not
      branch at the ends)
      \item replacing \texttt{sqrtf(E)~<~R} by a squared compare
      \item narrowing a \texttt{double} accumulator to \texttt{float}
    \end{itemize}
    \item Domain formula: Rec.~601 float MAC to Q15.
  \end{itemize}

  \item \textbf{Winner selection} instantiates~\eqref{eq:winner} under
  \texttt{/O2}. It keeps the lowest median of the contract's benchmark,
  and replaces \(P_1\) only when a candidate is faster by more than that
  candidate's sample standard deviation. When the contract has no
  benchmark, the first admitted candidate is emitted without a timing
  comparison.
\end{enumerate}

\paragraph{Scope.}
\textsc{IntC} is a C instance of \textsc{IntentCompile} on extracts.
\begin{itemize}
  \item The contract is \texttt{\#define} parameters, a YAML record of
  \(\mathcal{I}\), \(\mathcal{B}\), and \(\mathcal{O}\), and CBMC harnesses,
  with \(V\) bounded CBMC. \textsc{CheckRealiz} decides satisfiability of
  \(\mathcal{I}\), and the implemented conjunct of \(\mathcal{B}\) is an
  optional compiled-\texttt{sizeof} ceiling on file-scope objects.
  Unbounded operation sequences would need an inductive argument.
  \item The rewrite catalogs are small and can be extended the same way
  the domain formula is registered. For example, a linked queue or
  reordering struct fields to cut padding can be added that way. This
  \textsc{IntC} does not compile general C or C++ programs. For example, C++
  atomics, application-scale floats, and storage-format rewrites like
  CRS-to-ELLPACK are outside it.
  \item Lowering is MSVC~\texttt{/O2} on an Intel Core i7-14700K under
  Windows, and ranking compares \(P_1\) and the admitted programs under
  that same pair. That pair is \(\mathcal{O}\) for \textsc{IntC}, and
  another compiler or machine is another \(\mathcal{O}\).
\end{itemize}

\section{Evaluation}
\label{sec:evaluation}

This section asks three questions of \textsc{IntC}, in the order of the
subsections that answer them.
\begin{enumerate}
  \item Does \textsc{IntC} work on code extracted from real-world
  projects?
  \item Do as-if compilation of \(P_1\), and the oracle, budget filter, and
  \textsc{TheoryLower} score inside \textsc{IntentCompile}, behave as
  claimed?
  \item Under the same lowering backend that compiled \(P_1\), does
  \textsc{IntC} emit a candidate that is cheaper on
  \(\mathcal{O}\) than \(P_1\), or the same program?
\end{enumerate}
A last subsection names misses of \textsc{IntC}.
Runs use the machine and compilers of
Section~\ref{sec:implementation}. Tables are emitted by
\texttt{python reproduce.py} in the artifact
(\S\ref{sec:data-availability}). Solver times vary across runs. The
admit, timeout, and counterexample outcomes, and the rewrite names, are the
claims.

\subsection{Extracts from Real-World Projects}
\label{sec:catalog_table}

Each row of Table~\ref{tab:catalog} is one run of \textsc{IntC}
on a contract that states slack a developer could grant. How each extract
was cut from the upstream tree is in the artifact.

\begin{table}
\caption{One \textsc{IntC} run per row. Columns are the extract,
the catalog rule that fired, the harness bound \(K\), CBMC on \(P_1\) and
on the admitted program, and how that program compares to \(P_1\).}
\label{tab:catalog}
\begin{center}
\scriptsize
\begin{tabular}{llllll}
\toprule
Case & Rewrite & $K$ & $V$ on $P_1$ & $V$ on winner & As-if \\
\midrule
DPDK ring & pow2-mask & 4 & unknown/120.0s & proved/2.0s & not as-if \\
DPDK pad & pow2-mask & 4 & unknown/120.0s & proved/0.6s & unexposed \\
ClickHouse & pow2-mask & 4 & unknown/180.2s & proved/63.7s & not as-if \\
Zephyr & pow2-shift & 1 & proved/0.5s & proved/0.7s & not as-if \\
gain & pow2-mul & 1 & proved/0.1s & proved/0.1s & not as-if \\
PARSEC & aos-to-soa & 6 & proved/0.6s & proved/0.6s & layout \\
Flecs & predicated-update & 10 & proved/0.1s & proved/0.1s & layout \\
MachSuite & halo-1d & 8 & proved/0.1s & proved/0.1s & layout \\
CMSIS-DSP & align-loop-bound:BLOCK & 12 & proved/3.3s & proved/7.0s & not as-if \\
OGRE & narrow-int-field & 6 & proved/0.1s & proved/0.1s & layout \\
Box2D & squared-compare & 1 & proved/13.7s & proved/0.2s & not as-if \\
Opus & narrow-accum & 6 & proved/0.7s & proved/12.0s & not as-if \\
OpenCV & luma-q15 & 6 & unknown/60.0s & proved/0.1s & not as-if \\
PARSEC fat & aos-to-soa & 6 & proved/0.6s & proved/0.7s & layout \\
\bottomrule
\end{tabular}
\end{center}

\end{table}

\textbf{Case.} The cell names a tagged upstream situation whose hot loop we
extracted, except gain, which is a multiply-to-shift written for this
catalog. The DPDK ring and pad baselines are the 1000-slot sketches of
Section~\ref{sec:overview_intent}, not the library sources. ClickHouse's
baseline uses 10007 buckets, a prime, not the upstream power of two.
PARSEC and PARSEC fat are the same Euler step on four particles, the
latter on a 64-byte record that includes four explicit pad fields.
CMSIS-DSP is a two-tap FIR.

\textbf{Rewrite.} The cell is the catalog rule that produced the admitted
program. \textsc{TheoryLower} contributes remainder-to-mask, division-to-shift,
and multiply-to-shift. \textsc{GeneratePool} contributes the representation
rewrites and Rec.~601. In table order:
\label{sec:gold_rewrites}
\begin{itemize}
  \item DPDK ring: remainder-to-mask. CEGIS selects 1024.
  \textsc{GeneratePool} also admits pad-and-mask into that pool at \(K=4\).
  \item DPDK pad: remainder-to-mask on \texttt{STORAGE}, with
  \texttt{CAPACITY} still 1000. Fullness uses
  \(\texttt{tail}-\texttt{head}\) against \texttt{CAPACITY}, so enqueue
  fails at 1000 live packets. The extra 24 cells are a budget cost.
  \textsc{TheoryLower} and CEGIS produce that mask on the already-named
  \texttt{STORAGE}.
  \label{sec:unexposed_eval}
  \item ClickHouse: remainder-to-mask. \(\mathcal{I}\) requires at least
  10007 buckets. CEGIS selects 16384, the admissible power of two nearest
  that floor.
  \item Zephyr: \(\texttt{tick\_us / DIVISOR}\) to a shift. Relative error
  at most \(2.5\%\) versus \(\texttt{tick}/1000\) on 32-bit ticks at least
  640000, because integer division by 1024 overshoots that bound near
  600000. That floor is part of \(\mathit{Pre}\). CEGIS selects
  \(\texttt{DIVISOR}\) 1024 and \(\texttt{SHIFT}\) 10.
  \label{sec:decimation_eval}
  \item Gain: \(\texttt{sample * GAIN}\) with \(\texttt{GAIN} \in [8,16]\)
  to a left shift by 3. The harness asserts equality with
  \(\texttt{sample * GAIN}\) on 16-bit samples, so the admit is not a
  multiply by the baseline's 10.
  \label{sec:scale_eval}
  \item PARSEC: array-of-structs to structure-of-arrays.
  \label{sec:soa_eval}
  \item Flecs: predicated update on a 0/1 live flag.
  \item MachSuite: 1-D halo pad.
  \item CMSIS-DSP: loop bound aligned to a multiple of 4.
  \(\mathcal{I}\) is \(\texttt{BLOCK} \ge 5\) with search box
  \(\texttt{BLOCK} \le 64\); CEGIS selects 8.
  \item OGRE: 32-bit space tag narrowed to \texttt{uint8\_t}.
  \item Box2D: \(\texttt{sqrtf(d2) < R}\) to \(\texttt{d2 < R*R}\).
  \(\mathit{Post}\) is the integer disk test \(dx^2+dy^2 < R^2\); the
  \(\texttt{sqrtf}\) in \(P_1\) is not that test. The timed bench is
  \(10^6\) float pairs, outside \(\models_K\).
  \item Opus: \(\texttt{double}\) accumulator narrowed to \texttt{float}.
  The harness uses \(\mathit{LEN}=4\) operands in \(\{0,1,2,3\}\). The
  timed bench is \(L=256\) fractional floats, outside \(\models_K\).
  \item OpenCV: Rec.~601 float coefficients to a Q15 shift. The harness
  is that Q15 formula within one gray level, not the float Rec.~601 sum.
  \(\mathit{NPX}=4\) in the harness; the timed bench is \(1920\times 1080\),
  outside \(\models_K\).
  \item PARSEC fat: array-of-structs to structure-of-arrays.
\end{itemize}

{\bfseries\boldmath\(K\).} The cell is the harness bound for \(\models_K\): how many
operations or loop unwindings CBMC covers on that row
(\S\ref{sec:verifier}). Ring, pad, and ClickHouse use four operations from
an empty container. Zephyr, gain, and Box2D use one conversion, sample, or
compare. The rest are loop unwindings of the extract: 6 on PARSEC's four particles
and on OGRE, Opus, OpenCV, and PARSEC fat; 8 on MachSuite; 10 on Flecs;
12 on the two-tap FIR.

{\bfseries\boldmath\(V\) on \(P_1\) and on the winner.} Those cells are that CBMC
query. \(\mathsf{unknown}\) is a timeout. PARSEC and PARSEC fat prove the
Euler step at \(n=4\). CMSIS-DSP proves the two-tap filter. When both
programs prove, a longer time on the rewrite is the cost of that query.
Table~\ref{tab:catalog}'s Opus and CMSIS-DSP rows are the two where both
prove and the rewrite is the slower query.

\textbf{As-if.} The cell classifies the admitted program against \(P_1\).
\emph{not as-if} means a named parameter or a precision choice moved.
\emph{unexposed} means the backing store grew and caller-facing capacity
did not. \emph{layout} means a representation rewrite the harness treats as
unobservable.

\subsection{Behavior as Claimed}
\label{sec:microbench}

This subsection answers as-if compilation of \(P_1\), then the oracle, the
budget filter, and the \textsc{TheoryLower} score, in that order.

\textbf{As-if compilation of \(P_1\).}
The DPDK ring maps a sequence number \(i\) to a slot by remainder, and in
\(P_1\) that remainder is \(i \bmod 1000\). MSVC~\texttt{/O2} still emits a
reciprocal multiply for \texttt{i \% 1000} and \texttt{and ecx, 1023} for
\texttt{i \% 1024}, and it does not change 1000 to 1024. A microbenchmark of
that slot expression alone uses 15 interleaved trials of
\(2\cdot 10^8\) iterations and two seeds. A latency-bound chain measures
2.67\,ns versus 0.39\,ns, a factor of 6.8, and independent operations
measure 0.52\,ns versus 0.17\,ns, a factor of 3.1. That upper bound has no
queue memory traffic.

\label{sec:ablations}
Table~\ref{tab:ablations} is one switch at a time on that ring's contract
for the oracle and the budget.

\begin{table}
\caption{Four switches on the DPDK ring: FIFO oracle, \(\models_K\),
occupancy-only \(\mathit{Post}\), and an 8016-byte ceiling.}
\label{tab:ablations}
\begin{center}
\scriptsize
\begin{tabular}{p{0.34\linewidth}p{0.58\linewidth}}
\toprule
Switch & Outcome \\
\midrule
Contract oracle (FIFO harness) & proved/0.6s \\
Admission is $\models_K$ & K=4 proved/0.7s; K=10 counterexample/17.3s \\
Occupancy-only vs FIFO $\mathit{Post}$ & capacity-only proved/0.5s; FIFO counterexample/0.4s \\
BudgetNormalize at 8016\,B & emit P1 \\
\bottomrule
\end{tabular}
\end{center}

\end{table}

\textbf{Oracle.} The FIFO harness proves the admitted mask in 0.6\,s.
Admission is \(\models_K\): the colliding mask \(\texttt{\& 999}\) is
proved at \(K=4\) and refuted at \(K=10\), because slots collide once nine
or more elements are live.
\label{sec:bounded_gap}
An occupancy-only \(\mathit{Post}\) proves a LIFO stack that the FIFO
harness refutes.

\textbf{Budget filter.}
\label{sec:budget_eval}
\textsc{BudgetNormalize} at 8016 bytes, the 1000-slot baseline: CEGIS still
proves the 1024-slot mask at 8208 bytes. Algorithm~\ref{alg:budgetnormalize}
drops it and the emit is \(P_1\). The 1000-slot remainder still times out
at that \(K\).

\textbf{\textsc{TheoryLower} score.}
\label{sec:score_eval}
Forty remainder-site variants, on a ring at \(K=4\) and a hash set at
\(K=3\), have Spearman's \(\rho = 0.11\) between that score
(\S\ref{sec:implementation}) and CBMC time at a 120\,s timeout.
Every timeout still has a non-power-of-two remainder in the program, and
proved times are flat, so a score at or below the threshold is not a
decision that \(V\) will finish. A score that predicted CBMC time is a
different score (\S\ref{sec:discussion}). When TheoryLower finds no
remainder or division hole, it plants no seed: PARSEC's Euler step is that
case, and \textsc{GeneratePool} still admits structure-of-arrays from
\(P_1\).

\subsection{Cheaper on \(\mathcal{O}\), or \(P_1\)}
\label{sec:ring_eval}

This subsection answers whether ranking under the same lowering backend
emits a candidate cheaper on \(\mathcal{O}\) than \(P_1\), or \(P_1\).
Table~\ref{tab:timed} is the seven contracts that name a \texttt{bench.c}.
Min, max, and quartiles are in the artifact JSON.

\begin{table}
\caption{Contracts that name a \texttt{bench.c}, ordered by emit ratio,
\(P_1\) last. Columns are the extract, \(K\) on the timed harness, CBMC on
the admitted program, median wall-clock ns per unit of work of \(P_1\) and
of the admit under MSVC~\texttt{/O2}, and emit.}
\label{tab:timed}
\begin{center}
\scriptsize
\begin{tabular}{llllll}
\toprule
Case & $K$ & $V$ on admitted & $P_1$ ns & admitted ns & emit \\
\midrule
Box2D & 1 & proved/0.3s & $0.702 \pm 0.024$ & $0.108 \pm 0.015$ & $6.5\times$ \\
PARSEC fat & 6 & proved/0.6s & $1.203 \pm 0.039$ & $0.346 \pm 0.012$ & $3.5\times$ \\
ClickHouse & 4 & proved/62.4s & $9.455 \pm 0.097$ & $4.000 \pm 0.043$ & $2.4\times$ \\
DPDK ring & 10 & proved/47.8s & $1.128 \pm 0.017$ & $0.487 \pm 0.008$ & $2.3\times$ \\
Opus & 6 & proved/11.9s & $0.440 \pm 0.005$ & $0.273 \pm 0.003$ & $1.6\times$ \\
OpenCV & 6 & proved/0.1s & $0.840 \pm 0.011$ & $0.560 \pm 0.014$ & $1.5\times$ \\
CMSIS-DSP & 12 & proved/6.7s & $0.460 \pm 0.016$ & $0.927 \pm 0.062$ & $P_1$ \\
\bottomrule
\end{tabular}
\end{center}

\end{table}

\textbf{Case.} The cell is the Table~\ref{tab:catalog} extract whose timed
contract names that driver.

{\bfseries\boldmath\(K\).} The cell is the timed harness bound. Six of the
seven match Table~\ref{tab:catalog}. The DPDK ring is \(K=10\) here and
\(K=4\) in that table. Table~\ref{tab:catalog} used four operations from
an empty queue. This row uses ten, the bound that refutes
\texttt{\& 999} in Section~\ref{sec:bounded_gap}.

{\bfseries\boldmath\(V\) on admitted.} The cell is CBMC of the admitted
rewrite on that timed harness, including CMSIS-DSP where emit is \(P_1\)
and the cell is still the slower \(\texttt{BLOCK}=8\) admit. The ring row
is that query at \(K=10\). Table~\ref{tab:catalog}'s ring cell is the same
rewrite at \(K=4\). Solver seconds move across runs; the proved/unknown
outcome is the claim.

{\bfseries\boldmath\(P_1\) ns and admitted ns.} Those two cells are
\(\mathcal{O}\): wall-clock nanoseconds per unit of work on that
\texttt{bench.c} (pair, update, lookup, enqueue, mac, pixel, or FIR
block), lower is better. The stage compiles both programs with
MSVC~\texttt{/O2} and takes the median of 31 trials after warmup. A later
candidate replaces the current best only if it is faster by more than its
sample standard deviation.

\textbf{Emit.} The cell is the ratio of those medians, or \(P_1\) when the
admit is not faster by more than its sample standard deviation.
In table order:
\begin{itemize}
  \item Box2D: in-range pairs. MSVC~\texttt{/O2} still emits
  \(\texttt{sqrtss}\) from the baseline, so the squared compare is the
  ranking win. The listing is in the artifact.
  \item PARSEC fat: the Euler step at \(N=100000\). The catalog query
  still proves \(n=4\).
  \item ClickHouse: fill 4000 keys, then \texttt{contains}. The ratio
  mixes a larger table and a lower load factor with the mask.
  \label{sec:hash_eval}
  \item DPDK ring: a dependent enqueue/dequeue chain through a partly
  full queue. The two programs disagree on the 1001st consecutive
  enqueue. The isolated slot chain of
  Section~\ref{sec:microbench} overstates that whole-queue gain. The
  baseline remainder times out at 300\,s for \(K=4,6,8\) in the recorded
  sweep, and at 120\,s in Table~\ref{tab:catalog}.
  \item Opus: the inner-product kernel.
  \item OpenCV: luma pixels.
  \item CMSIS-DSP: one \(\texttt{filter()}\) call. Extra trips from
  rounding 5 to 8 are slower, so this row emits \(P_1\).
  Table~\ref{tab:catalog} used a contract with no \texttt{bench.c}, so
  that row emitted the admitted \(\texttt{BLOCK}=8\) rewrite without a
  timing comparison (\S\ref{sec:implementation}).
\end{itemize}

\subsection{Misses of This Instance}
\label{sec:failures}

This subsection names where the three answers fail on \textsc{IntC}.
Table~\ref{tab:failures} is those misses.

\begin{table}
\caption{Misses of \textsc{IntC}: never generated, budget, slower
admit, weak harness, bound too small, driver, and unnamed cost.}
\label{tab:failures}
\begin{center}
\scriptsize
\begin{tabular}{llp{0.42\columnwidth}}
\toprule
Kind & Instance & What happens \\
\midrule
Never generated & linked queue, CRS-to-ELLPACK, swap-and-pop & not in the catalog \\
Budget & 1024-slot mask over 8016\,B & dropped, emit \(P_1\) \\
Slower admit & FIR \(\texttt{BLOCK}\) 5\(\to\)8 & extra trips, emit \(P_1\) \\
Weak harness & occupancy-only LIFO & FIFO harness refutes it \\
Bound too small & \(\texttt{\& 999}\) & proved at \(K=4\), refuted at \(K=10\) \\
Driver & slot chain vs ring \texttt{bench.c} & \(6.8\times\) vs \(2.3\times\), same rewrite \\
Unnamed cost & false sharing & not measured \\
\bottomrule
\end{tabular}
\end{center}
\end{table}

\textbf{Kind.} The cell is how the miss happens.

\textbf{Instance.} The cell is the extract or rewrite.

\textbf{What happens.} The cell is the outcome. In table order:
\begin{itemize}
  \item Never generated: a linked queue, a CRS-to-ELLPACK layout, and a
  swap-and-pop compaction are not catalog sites, so they never become
  candidates. A CRS-to-ELLPACK rewrite that kept the stored values would
  still stream the extra zeros ELLPACK pads.
  \item Budget: the 1024-slot mask proves and is 8208 bytes. The 8016-byte
  ceiling drops it and the emit is \(P_1\)
  (Section~\ref{sec:budget_eval}).
  \item Slower admit: FIR \(\texttt{BLOCK}\) 8 does extra trips on one
  \(\texttt{filter()}\) call, so ranking emits \(P_1\)
  (Section~\ref{sec:ring_eval}).
  \item Weak harness: an occupancy-only \(\mathit{Post}\) proves a LIFO
  stack that the FIFO harness refutes
  (Section~\ref{sec:microbench}).
  \item Bound too small: \(\texttt{\& 999}\) proves at \(K=4\) and is
  refuted at \(K=10\) (Section~\ref{sec:bounded_gap}).
  \item Driver: the isolated slot chain is \(6.8\times\). The ring
  \texttt{bench.c} is \(2.3\times\) on the same rewrite
  (Section~\ref{sec:microbench}).
  \item Unnamed cost: concurrent false sharing is not in any named
  driver, so it cannot reject a candidate.
\end{itemize}

\section{Related Work}
\label{sec:related}

\paragraph{Compiling against a Specification.}
Leroy distinguishes compilers that preserve a program's behavior from
compilers that preserve only a specification, and notes that the latter
leave more room for optimization~\cite[\S2.1]{leroy2009backend}. Refinement
calculi make any program that refines a specification a correct
implementation~\cite{back1998refinement,morgan1994programming}. Fiat
refines specifications with deliberate freedom into proved implementations
in Coq~\cite{delaware2015fiat}. Data representation synthesis, Cozy, and
Primrose choose among the implementations a specification
admits~\cite{hawkins2011representation,loncaric2016cozy,Qin2023Primrose}.

Those systems compile a specification: Fiat and Cozy in a language built
for that compilation, Primrose from annotations on Rust into a library
container. \textsc{IntentCompile} uses a general contract with a
general-purpose program the developer wrote as the seed of the search.
The compiler treats that seed as a proxy of the developer's mental model of
the problem and searches from it for a similar but cheaper solution.
Section~\ref{sec:ideal} states the problem for a general-purpose language
of the kind developers already run in production. \textsc{IntC} is the C
instance this paper evaluates. Table~\ref{tab:related} places that instance
against specification compilers, schedulers, and as-if-relaxing tools.

\begin{table}
\caption{Columns are the host language, whether \(\theta\) is a named
parameter the developer writes, whether the search includes a rewrite
as-if would reject, the admission oracle, and whether the compiler can
emit \(P_1\) when the search fails. \textsc{IntentCompile} is the
algorithm. \textsc{IntC} is the C instance.}
\label{tab:related}
\centering
\scriptsize
\begin{tabular}{lccccc}
\toprule
System & Language & Named \(\theta\) & Non-as-if slack & Sound \(V\) & Fallback \(P_1\) \\
\midrule
Fiat~\cite{delaware2015fiat} & Coq & in the spec & yes & Coq & no \\
Cozy~\cite{loncaric2016cozy} & spec to Java & in the spec & yes & synthesis checker & no \\
Primrose~\cite{Qin2023Primrose} & annotated Rust & container spec & yes & SMT & library pick \\
STOKE~\cite{schkufza2014fp} & x86 asm & error bound & numeric & validator & best-so-far \\
Alive2~\cite{lopes2021alive2} & LLVM IR & no & no & yes & n/a \\
Halide/Exo~\cite{ragan2013halide,ikarashi2022exo} & own language & schedule & no & lowering & n/a \\
PetaBricks~\cite{ansel2009petabricks} & own language & autotune knobs & authored menu & no & no \\
\textsc{IntC} & C & \(\mathcal{I}\) & yes & \(\models_K\) & yes \\
\bottomrule
\end{tabular}
\end{table}

\paragraph{Scheduling Languages.}
Halide and Exo keep a DSL algorithm, including sizes, fixed and search its
schedule~\cite{ragan2013halide,ikarashi2022exo}. \textsc{IntentCompile}
searches structured refactors of a general-purpose program, starting from
the current implementation and a contract that names which of its choices
are slack.

\paragraph{Optimizing under Relaxed Correctness.}
Production compilers already ship named, opt-in relaxations of as-if.
\texttt{-ffast-math} is one example, a fixed license to reassociate and
otherwise rewrite IEEE-754 arithmetic, and scoped floating-point pragmas
are another. STOKE optimizes floating-point code against
a user-chosen error bound~\cite{schkufza2014fp}. Sound compilation of reals
and Chisel choose
any implementation within a stated
tolerance~\cite{darulova2014sound,misailovic2014chisel}. Conditionally
correct superoptimization exploits facts about inputs beyond undefined
behavior~\cite{sharma2015conditionally}. Carbin et al.\ verify acceptability
properties of relaxed programs~\cite{carbin2012relaxed}. Loop perforation
trades accuracy for speed under empirical
checks~\cite{sidiroglou2011managing}. PetaBricks lets the developer write
several algorithms for one computation and autotunes among them by
measurement~\cite{ansel2009petabricks}.

They each implement a particular relaxation, an opt-in compiler flag, an
error bound, extra input facts, perforation, or an authored menu. Because \textsc{IntentCompile}
does not take as-if as the bound, a contract can name several of those
freedoms together, and the search is a structured refactor of the current
implementation.

\paragraph{Translation Validation and Superoptimization.}
Translation validation~\cite{pnueli1998translation,necula2000translation}
and CompCert~\cite{leroy2009backend} certify as-if: CompCert proves that the
backend preserves the source's observable behavior. Alive2 checks that a
transformation refines its LLVM
source~\cite{lopes2021alive2, lee2017taming}. Souper synthesizes peephole
optimizations that the same check
admits~\cite{sasnauskas2017souper}. Recent systems pair language models
with these checkers~\cite{taneja2025llmvectorizer,fang2026llmveriopt}.
Rinard finds that agents under relational equivalence can pick slower
algorithms because they are easier to prove
equivalent~\cite{rinard2026quantitative}.

Those checkers ask whether a candidate is equivalent to the source or
refines it. \textsc{IntentCompile} keeps the same split of untrusted
proposer and sound checker, and the checker asks whether the candidate
satisfies the contract, which is what admits a 1024-cell ring from a
1000-cell baseline.

\paragraph{Lifting and Verification-Friendly Programs.}
Verified lifting translates code into a higher-level form that can be
verified and then regenerated~\cite{kamil2016stng,bhatia2024llmlift}, with
the lifted specification proved equivalent to the source. Wagner et al.\
compile a separate, verification-friendly build~\cite{wagner2013overify}.
Instrumentation can make programs easier for automated tools to
verify~\cite{amilon2023instrumentation}. Ivy restructures protocols into
decidable fragments by
hand~\cite{padon2016ivy,padon2017paxosepr}. Predicting solver cost is
studied for choosing among tools~\cite{healy2017where4,leeson2022graves}.

A specification proved equivalent to the source keeps the legal set inside
as-if of that source. \textsc{IntentCompile} generates candidates from the
current implementation and the contract. \textsc{TheoryLower} rewrites
toward a sketch \(V\) can decide, and \(V\) admits against the contract.
A measurement of \(\mathcal{O}\) may not need that proof, but the compiler
still ranks only among admitted programs so a cheaper but illegal candidate
is not emitted.

\paragraph{Data Layout in Compilers.}
Production compilers restructure data when they can prove the change
invisible~\cite{zhao2007forma,hundt2006practical,lattner2005pool}. MEMOIR
makes collections explicit in IR~\cite{mcmichen2024memoir}. Annotation-based
AoS-to-SoA conversions are trusted rather than
checked~\cite{radtke2025annotation}.

Those passes are unexposed slack (\S\ref{sec:slack}): extra cells or a
different layout the interface does not observe. \textsc{IntentCompile}
can name that slack and logical-parameter slack, a capacity the caller can
observe, in one contract, and \(V\) admits the refactor.

\paragraph{Objectives.}
Objectives that rank programs come in many kinds. Autotuners such as
PetaBricks rank by measured
runtime~\cite{ansel2009petabricks}. McCabe's cyclomatic
complexity~\cite{mccabe1976complexity} and Martin's coupling and
abstraction~\cite{martin2000design} are computed from the program text.
\(\mathcal{O}\) in Section~\ref{sec:opt_contract} is that supplied ranking
over programs the contract admits. \textsc{IntC} fills it with a
profile-guided measurement, wall-clock time on a named driver.

\section{Discussion and Future Work}
\label{sec:discussion}

\paragraph{Writing \(\Phi\).}
Decoupling search from \(P_1\) moves the specification burden onto the
developer. That cost has limited adoption of contract-based tools.
A contract inferred from the baseline tends to freeze incidental choices,
such as a capacity of exactly 1000, and removes the slack the problem
depends on. Contracts have to come from intent, not from the code.
Translating a goal-oriented requirements model into \(\mathit{Pre}\),
\(\mathit{Post}\), and \(\mathcal{I}\) is future work. One route is a
KAOS refinement of a goal tree down to atomic requirements in a formal
specification language~\cite{vanlamsweerde2001gore}.

\paragraph{Under-Specification.}
The legal set is the slack in \(\Phi_{\text{func}}\), including slack the
developer did not intend. A FIFO ring whose contract omits order admits a
faster queue that reorders packets, and \(V\) will prove it. Checking a
proposed contract against the baseline catches contracts that are too
strong, not contracts that are too weak~\cite{misu2026specharness}. A missing
requirement is a proved program the developer did not want. When a
deployment or a rejected candidate surfaces a gap, that requirement
should become a contract clause that later compilations re-check, so the
same regression cannot return unless the requirements change. Completing
\(\Phi\) that way is the developer's cycle around the compiler, not a
pass of Section~\ref{sec:compiler}. The same holds for a cost that is
missing from \(\mathcal{B}\) or \(\mathcal{O}\): concurrent false sharing
on a padded ring cannot reject a candidate until it is named. Closing
gaps permanently is the central open problem for intent compilation.

\paragraph{Temporarily Relaxing Constraints.}
Loosening \(\mathit{Pre}\) or \(\mathit{Post}\) enlarges
\(\mathrm{Sat}(\Phi_{\text{func}})\), so the search can admit a program the
tight contract would reject. Putting those clauses back makes that program
illegal, and \(V\) has to run again. Loosening \(\mathcal{B}\) does not
change legality: a candidate admitted under a looser memory ceiling can
still satisfy \(\Phi_{\text{func}}\) after the tight ceiling is restored,
and no observable test fails. The compiler also bounds how much of
\(\mathrm{Gen}\) it explores and how long \(V\) may run; loosening those
changes which candidates finish in a compile, not
\(\mathrm{Sat}(\Phi_{\text{func}})\). Restoring a temporary relaxation is
future work: re-running \(V\), re-checking \(\mathcal{B}\), or re-searching
under the tight bounds.

\paragraph{Unbounded Verification.}
\textsc{IntC}'s claim is \(\models_K\). An unbounded proof of a stateful
component is the invariant-and-abstraction obligation Ivy discharges by
restructuring into a decidable
fragment~\cite{padon2016ivy,padon2017paxosepr}. Carrying that obligation
on a candidate, rather than stopping at \(K\), is future work.

\paragraph{Where to Spend Compilation Time.}
\textsc{IntC} runs \(V\) before it measures \(\mathcal{O}\), so a
candidate can consume a timeout before a cheap profile would have
dropped it. Other schedules over the same \(\mathrm{Gen}\) include
estimating \(\mathcal{O}\) first, running a cheap \(\models_K\) before a
longer query, and keeping candidates and proofs as artifacts so a later
compile reuses them.

\paragraph{A Score that Predicts \(V\).}
\textsc{IntC}'s TheoryLower score orders descent, and against CBMC time it
has Spearman's \(\rho = 0.11\). A score that prefers the rewrite expected
to shorten the query has to be fit to this \(V\), and predicting SMT and
verifier runtime from goal or program structure is already
studied~\cite{healy2017where4,leeson2022graves}.

\paragraph{Language Models as Compiler Inputs.}
\label{sec:discussion_lm}
A model can be the untrusted \(\mathrm{Gen}\) of
Theorem~\ref{thm:restricted} at compile
time~\cite{taneja2025llmvectorizer,fang2026llmveriopt,jiang2026segabench},
or, for a deterministic compile, help write the catalogs, TheoryLower
scores, and benchmarks an instance enumerates. It can also help with the
specification burden of writing \(\Phi\) by drafting
\(\Phi_{\text{func}}\) from a prose
requirement~\cite{ma2025specgen,ugare2026fun2spec} as a candidate for
\(V\), and by helping write the baseline \(P_1\). \textsc{IntC} is the prototype
for a C instance; a model can help instantiate those inputs for another
language.

\paragraph{Language Models as Developers.}
Professional developers commonly use AI tools on coding tasks: 51\%
report daily use in the 2025 Stack Overflow
survey~\cite{stackoverflow2025ai}, including multi-turn conversations on
real GitHub work~\cite{hao2024chatgpt}. Later rewrites in that setting
often break requirements that earlier turns had already
met~\cite{huang2026regression,yan2026slump}.
\textsc{IntentCompile} takes those requirements as a checkable
\(\Phi\) and the program as \(P_1\), which a model can
fill~\cite{lahiri2026intent}. The next rewrite is then a compile against
\(\Phi\). Named requirements change when intent changes, which is less
often than the program is rewritten, so a human can maintain \(\Phi\)
while the model rewrites \(P\).

\section{Conclusion}
\label{sec:conclusion}

An as-if compiler preserves the source, including incidental choices the
requirements do not fix, so it may omit a cheaper legal program.
\textsc{IntentCompile} takes a baseline and a checkable contract,
generates candidates, and admits those a sound verifier proves. It ranks
the admitted set by \(\mathcal{O}\) under the same lowering backend that
compiled \(P_1\), and emits \(P_1\) when no admitted candidate is cheaper.
\textsc{IntC} is a prototype C instance of that restricted compiler.

On the motivating example, \textsc{IntC} admits the 1024-cell ring that
MSVC~\texttt{/O2} will not emit from a 1000-cell remainder. On extracts
whose contracts name a benchmark, the emit is cheaper on \(\mathcal{O}\)
than \(P_1\) or is \(P_1\). The catalogs, the bound \(K\), and writing
\(\Phi\) from intent are the limits of \textsc{IntC}.


\section*{Data-Availability Statement}
\label{sec:data-availability}

An artifact accompanies this paper: \textsc{IntC}, a Python source-to-source
compiler for C, the contracts and C extracts of
Section~\ref{sec:evaluation}, CBMC harnesses, and the script that rebuilds
the catalog, ablation, and timed-contract tables. Admission is
\(\models_K\) with CBMC 6.11.0. Measurement and listings use
MSVC 19.29~\texttt{/O2} on Windows. From \texttt{poc/}, after
\texttt{pip install -e .}, \texttt{python reproduce.py} rebuilds those
tables. Solver times are not bit-stable. Admit, timeout, and
counterexample outcomes, and rewrite names, must match
\texttt{ARTIFACT\_README.md}.
The artifact is \url{https://github.com/LightShield/IntC}.

\bibliographystyle{ACM-Reference-Format}
\bibliography{references}

\appendix
\section{Proofs}
\label{app:proofs}

Throughout, ``result'' means the observable part of
\(\llbracket P \rrbracket(\sigma)\), and \(\mathit{Acc}_\theta(\sigma)\) is
the set of observable results that \(\mathit{Post}_\theta\) accepts on
\(\sigma\).

\paragraph{Proposition~\ref{prop:refine} (As-If Programs Remain Legal).}
If \(P \models \Phi_{\text{func}}\), \(P' \approx P\), and
\(\theta_{P'} = \theta_P\), then \(P' \models \Phi_{\text{func}}\).
\begin{proof}
Suppose \(P' \not\models \Phi_{\text{func}}\).
Then either \(\theta_{P'} \notin \mathcal{I}\), or some
\(\sigma \in \mathit{Pre}\) is a non-terminating, undefined, or unacceptable
run of \(P'\). The first case is impossible because
\(\theta_{P'} = \theta_P \in \mathcal{I}\). In the second case,
\(P' \approx P\) gives the same observable result for \(P\) on \(\sigma\),
so \(P\) also violates the contract, a contradiction.
\end{proof}

\paragraph{Proposition~\ref{prop:expansion} (Legal Set Expands Exactly on Observable Slack).}
Assume local patching, and that each \(\mathit{Post}_\theta\) with
\(\theta \in \mathcal{I}\) is realized by some program. Let
\(P \models \Phi_{\text{func}}\). Then
\(\mathrm{Sat}(\Phi_{\text{func}})\) contains a program that is not
observationally equivalent to \(P\) iff the contract has input, output, or
logical-parameter slack, counting only results some program produces.
Unexposed slack is not in that list: it need not change those results.
\begin{proof}
Assume also that all results below are produced by some program.
Let \(o_P(\sigma)\) be \(P\)'s result on \(\sigma\).

(\(\Leftarrow\))
\emph{Input slack.} Let \(\sigma_0 \notin \mathit{Pre}\) and let \(o\) be a
producible result different from \(o_P(\sigma_0)\). Patch \(P\) at
\(\sigma_0\). The patch agrees with \(P\) on \(\mathit{Pre}\), so it
satisfies the contract, and differs from \(P\) on \(\sigma_0\).

\emph{Output slack.} Let \(\sigma_0 \in \mathit{Pre}\) and
\(\theta = \theta_P\) with two accepted results, one of them \(o \neq
o_P(\sigma_0)\). (If the slack is only at another parameter, it is the
logical-parameter case.) Patch \(P\) at \(\sigma_0\) to produce \(o\).

\emph{Logical-parameter slack.} Let \(\sigma_0 \in \mathit{Pre}\) and
\(\theta_1, \theta_2 \in \mathcal{I}\) with
\(\mathit{Acc}_{\theta_1}(\sigma_0) \neq \mathit{Acc}_{\theta_2}(\sigma_0)\).
Both sets are nonempty. Some accepted result \(o\) at one of those
parameters differs from \(o_P(\sigma_0)\). Start from a program for that
parameter and patch it at \(\sigma_0\).

(\(\Rightarrow\)) If there is no such slack, then \(\mathit{Pre} = \Sigma\),
each \(\mathit{Acc}_\theta(\sigma)\) has at most one producible element, and
those elements coincide for all \(\theta \in \mathcal{I}\). That element is
\(o_P(\sigma)\). Any legal \(P'\) must produce it everywhere, so
\(P' \approx P\).
\end{proof}

Unexposed slack is omitted because it need not change observable API results.

\paragraph{Lemma~\ref{lem:lowering} (Lowering Preserves a Proved Contract).}
If \(P^* \models \Phi_{\text{func}}\) (respectively \(\models_K\)), then
\(\mathcal{L}_{\text{Target}}(P^*) \models \Phi_{\text{func}}\)
(respectively \(\models_K\)).
\begin{proof}
Let \(B = \mathcal{L}_{\text{Target}}(P^*)\). Parameters of \(B\) are those
of the source given to the backend. If \(P^*\) satisfies \(\Phi_{\text{func}}\)
(or \(\models_K\)) then it has no undefined behavior on the covered inputs
(or traces), terminates with an accepted result, and the compiler
correctness assumption copies that result to \(B\).
Proposition~\ref{prop:refine} is the unbounded case.
\end{proof}

\paragraph{Theorem~\ref{thm:restricted} (Soundness of a Restricted Intent Compiler).}
Let \(V\) be sound for \(\models\) or for \(\models_K\), and let the compiler
emit source \(P^*\) and binary
\(\mathrm{Bin}^* = \mathcal{L}_{\text{Target}}(P^*)\).
\begin{enumerate}
  \item If \(P^* \neq P_1\), then \(P^*\) is in \(\mathrm{Gen}(P_1, R, D)\),
  \(V\) proved it, and it passed \(\mathcal{B}\). Hence
  \(P^* \models \Phi_{\text{func}}\) or \(P^* \models_K \Phi_{\text{func}}\),
  matching \(V\), and the same judgment holds of \(\mathrm{Bin}^*\) by
  Lemma~\ref{lem:lowering}.
  \item The measured cost of \(\mathrm{Bin}^*\) is at most that of
  \(\mathcal{L}_{\text{Target}}(P_1)\), when measurement runs.
  \item If \(P_1 \models \Phi_{\text{func}}\) (or \(\models_K\)), then
  \(\mathrm{Bin}^*\) satisfies the same judgment in every case.
\end{enumerate}
\begin{proof}
Winner selection (Algorithm~\ref{alg:winner}) returns \(P_1\) or a member of \(\mathcal{A}\). If
\(P^* \neq P_1\), then \(P^* \in \mathcal{A}\), so \(P^* \in \mathrm{Gen}\),
\(V\) answered \(\mathsf{proved}\), and \(\mathcal{B}\) passed.
Soundness of \(V\) gives \(\models\) or \(\models_K\).
Lemma~\ref{lem:lowering} transfers the judgment to the binary. Measured
cost is minimized over a set containing
\(\mathcal{L}_{\text{Target}}(P_1)\). If \(P^* = P_1\) and \(P_1\) already
satisfies the judgment, apply Lemma~\ref{lem:lowering} to \(P_1\).
\end{proof}

\section{Pipeline Algorithms}
\label{app:algorithms}

This appendix gives each stage of \textsc{IntentCompile} as an algorithm:
inputs, outputs, the promises that algorithm makes, and a proof of those
promises. Every stage claims termination, relative to its oracles returning
(the solver for \(\mathcal{I}\), \(V\), the checker for \(\mathcal{B}\),
lowering, and measurement). Parameters such as the TheoryLower score, the
CEGIS distance, the attempt cap, the checker for \(\mathcal{B}\), and the
rewrite catalogs are inputs of the algorithms.
Section~\ref{sec:implementation} instantiates them. Hole domains are finite. \(V\) returns
\(\mathsf{proved}\), a counterexample, or \(\mathsf{unknown}\), and respects
its time budget. A rewrite in a catalog has a syntactic applicability
condition and optional side constraints on holes. Nothing below claims that
\(\mathrm{Sat}(\Phi_{\text{func}})\) is searchable, that the catalogs are
complete, or that \(V\) decides every query.

\subsection{\textsc{CheckRealiz}}

\begin{algo}[\textsc{CheckRealiz}]
\label{alg:checkrealiz}
Input: \(\Phi_{\text{func}} = \langle \mathit{Pre}, \mathit{Post},
\mathcal{I}\rangle\) and a solver for \(\mathcal{I}\) that answers
satisfiable, unsatisfiable, or \(\mathsf{unknown}\). A time budget on that
solver, when present, is mapped to \(\mathsf{unknown}\).
Output: \(\mathsf{continue}\), or a stop
(\(\mathsf{vacuous}\), \(\mathsf{unsat}\) with a core, or
\(\mathsf{unknown}\)).
\begin{enumerate}
  \item If \(\mathit{Pre} = \emptyset\), report \(\mathsf{vacuous}\) and
  stop.
  \item Ask the solver whether \(\mathcal{I}\) is satisfiable.
  \item If the answer is unsatisfiable, report \(\mathsf{unsat}\) with the
  solver's unsatisfiable core and stop.
  \item If the answer is \(\mathsf{unknown}\), stop.
  \item If the answer is satisfiable, report \(\mathsf{continue}\).
\end{enumerate}
\end{algo}

\begin{lemma}[\textsc{CheckRealiz}]
\label{lem:checkrealiz}
Assume the solver for \(\mathcal{I}\) is sound.
\begin{enumerate}
  \item \emph{Termination.} If the solver for \(\mathcal{I}\) returns in
  finite time, so does the algorithm. A timeout is \(\mathsf{unknown}\) and
  is a return.
  \item \emph{No candidate.} The algorithm emits no program.
  \item \emph{Vacuous stop.} If it reports \(\mathsf{vacuous}\), then
  \(\mathit{Pre} = \emptyset\), so every program satisfies
  \(\Phi_{\text{func}}\) and the contract states no intent.
  \item \emph{Unsat stop.} If it reports \(\mathsf{unsat}\), then
  \(\mathcal{I}\) is empty, so no program satisfies
  \(\Phi_{\text{func}}\).
  \item \emph{Unknown stop.} An \(\mathsf{unknown}\) answer is a stop, not a
  proof that \(\mathcal{I}\) is inhabited or empty.
  \item \emph{Continue.} A \(\mathsf{continue}\) result is only that
  \(\mathcal{I}\) has a model. It is not a candidate, and it does not decide
  whether \(\mathit{Post}\) is realizable on \(\mathit{Pre}\).
\end{enumerate}
\end{lemma}
\begin{proof}
Steps 1 and 3--5 are local. The only query is step 2. If the solver
returns satisfiable, unsatisfiable, or \(\mathsf{unknown}\), one of those
steps returns, which is (1). The algorithm has no step that constructs or
admits a program, so (2) holds. Step 1 returns \(\mathsf{vacuous}\) exactly
when \(\mathit{Pre} = \emptyset\). Satisfaction of \(\Phi_{\text{func}}\)
quantifies over \(\mathit{Pre}\), so an empty precondition is true of every
program, which is (3). Step 3 returns \(\mathsf{unsat}\) only on a sound
unsat answer, which means \(\mathcal{I}\) has no valuation. A program that
satisfies \(\Phi_{\text{func}}\) must have \(\theta_P \in \mathcal{I}\), so
no such program exists, which is (4). Step 4 returns on
\(\mathsf{unknown}\) without a model or a core, which is (5). The query in
step 2 mentions \(\mathcal{I}\) only. A model of \(\mathcal{I}\) does not
produce a program and does not prove that every \(\sigma \in \mathit{Pre}\)
has an accepted result under \(\mathit{Post}\), which is (6).
\end{proof}

The stronger query \(\forall \sigma \in \mathit{Pre}.\ \exists o.\
(\sigma, o) \in \mathit{Post}_\theta\) is not this stage. An unsatisfiable
harness is discovered later, when no candidate is proved. Termination is
relative to the solver: the stage does not decide an arbitrary theory of
\(\mathcal{I}\), and a solver with no budget may diverge.

\subsection{\textsc{TheoryLower}}

Let \(\mathit{score}\) map sketches to \(\delta\cdot\mathbb{N}\) for some
granularity \(\delta > 0\), and let \(\tau\) be a threshold at which
rewriting stops. Both are inputs. The intended reading is a weighted count
of constructs outside a chosen cheap set for the source language and
\(V\), so descent aims at that set. Any score with those values for which
every applied rewrite strictly decreases the score makes the lemma hold.
\textsc{IntC}'s score is the weighted feature count of
Section~\ref{sec:implementation}. \(\mathit{score}(S) \le \tau\) does not
imply that \(V\) decides \(S\).

\begin{algo}[\textsc{TheoryLower}]
\label{alg:theorylower}
Input: baseline \(P_1\), catalog \(R\), \(\mathit{score}\), and
threshold \(\tau\).
Output: a sketch \(S\) tagged \(\mathsf{below\text{-}threshold}\) or
\(\mathsf{stuck}\), or no sketch.
\begin{enumerate}
  \item Set \(S \gets P_1\). Rewrites, not this step, introduce holes.
  \item If \(\mathit{score}(S) \le \tau\), return no sketch, tagged
  \(\mathsf{below\text{-}threshold}\).
  \item Among pairs of a rewrite \(r \in R\) and a site of \(S\) where
  \(r\) applies, choose one that minimizes \(\mathit{score}(r(S))\). If no
  pair gives a strict decrease: if a rewrite has already been applied,
  return \((S, \mathsf{stuck})\); otherwise return no sketch, tagged
  \(\mathsf{stuck}\).
  \item Set \(S \gets r(S)\). If \(\mathit{score}(S) \le \tau\), return
  \((S, \mathsf{below\text{-}threshold})\). Otherwise go to step 3.
\end{enumerate}
\end{algo}

\begin{lemma}[\textsc{TheoryLower}]
\label{lem:theorylower}
\begin{enumerate}
  \item \emph{Termination.} The algorithm returns after at most
  \(\mathit{score}(P_1)/\delta\) rewrites.
  \item \emph{Catalog.} Every rewrite it applies is from \(R\), at a site
  that matches that rewrite. Holes come from those rewrites.
  \item \emph{Exits.} CEGIS receives a sketch only if a rewrite applied.
  That sketch has \(\mathit{score} \le \tau\), or no remaining rewrite in
  \(R\) lowers the score. If no rewrite applied, there is no TheoryLower
  seed.
  \item \emph{Unverified.} No returned sketch is submitted to \(V\) at this
  stage.
\end{enumerate}
\end{lemma}
\begin{proof}
Each applied rewrite lowers \(\mathit{score}\) by at least \(\delta\), and
the score is nonnegative, so that can happen at most
\(\mathit{score}(P_1)/\delta\) times, which is (1). Step 3 selects from
\(R\) at matching sites, and step 1 does not punch holes, which is (2).
Step 2 returns no sketch. Step 3 returns a sketch only if one was already
applied, or no sketch. Step 4 returns a rewritten sketch, which is (3).
No step calls \(V\), which is (4).
\end{proof}

The stage makes no semantic claim. The sketch may violate the contract for
every hole value, and later stages decide which fillings survive. The loop
is greedy: a rewrite that first raises the score and enables a larger drop
later is never taken. A \(\mathsf{stuck}\) result does not abort
\textsc{GeneratePool}, which starts from \(P_1\).

\subsection{\textsc{Baseline-Directed CEGIS}}

Let \(S[h]\) be a sketch with holes ranging over a finite domain \(D\), and
let \(\mathit{dist}\) be a distance on \(D\). Closeness is relative to
\(\mathit{dist}\). \textsc{IntC} minimizes distance of baseline holes to
\(\theta_{P_1}\) and extra holes subject to the side constraints
(Section~\ref{sec:implementation}).

\begin{algo}[\textsc{Baseline-Directed CEGIS}]
\label{alg:cegis}
Input: sketch \(S\), domain \(D\), side constraints, \(\mathcal{I}\),
baseline parameters \(\theta_{P_1}\), verifier \(V\) with budget \(t\), and
a finite attempt cap \(N\).
Output: a filled program \(S[h]\), or \(\mathsf{fail}\).
\begin{enumerate}
  \item Set \(R_{\mathrm{ex}} \gets \emptyset\) and \(n \gets 0\).
  \item If \(n = N\), return \(\mathsf{fail}\).
  \item Find \(h \in D \setminus R_{\mathrm{ex}}\) that meets the side
  constraints, gives \(\theta_{S[h]} \in \mathcal{I}\), and minimizes
  \(\mathit{dist}(h, \theta_{P_1})\). If none exists, return
  \(\mathsf{fail}\).
  \item Run \(V(S[h], \Phi_{\text{func}}, t)\).
  \item If the answer is \(\mathsf{proved}\), return \(S[h]\).
  \item If the answer is a counterexample, add \(h\) to
  \(R_{\mathrm{ex}}\). Do not add the failing input as a constraint on
  later proposals. Set \(n \gets n+1\) and go to step 2.
  \item If the answer is \(\mathsf{unknown}\), return \(\mathsf{fail}\).
\end{enumerate}
\end{algo}

\paragraph{Theorem~\ref{thm:cegis} (Properties of Baseline-Directed CEGIS).}
Let \(S[h]\) range over a finite hole domain \(D\), and let \(N\) be the
attempt cap (\(N = |D|\) if there is no cap).
\begin{enumerate}
  \item If CEGIS returns \(S[h]\), then \(V\) proved \(S[h]\).
  \item \emph{Termination.} If \(V\) returns on each query, CEGIS returns
  after at most \(\min(N, |D|)\) iterations. Step 3 searches a finite set,
  so it returns.
  \item If some \(h^\circ \in D\) makes \(S[h^\circ]\) satisfy the contract
  (or \(\models_K\)), \(V\) never answers \(\mathsf{unknown}\), and there is
  no attempt cap, then CEGIS returns a satisfying assignment closest to
  \(\theta_{P_1}\).
\end{enumerate}
\begin{proof}
(1) The only success return is step 5, which requires
\(\mathsf{proved}\).

(2) \(V\) is assumed to return (\(\mathsf{proved}\), a counterexample, or
\(\mathsf{unknown}\)) because it respects \(t\). Each unsuccessful
iteration that does not fail on \(\mathsf{unknown}\) adds one \(h\) to
\(R_{\mathrm{ex}}\) and never retries it. The loop also stops when
\(n = N\). Together that is at most \(\min(N, |D|)\) iterations. Step 3
enumerates a finite remainder of \(D\).

(3) Let \(H^\circ = \{h' \in D \mid S[h'] \models \Phi_{\text{func}}\}\)
(or \(\models_K\)), nonempty by assumption. Soundness of \(V\) keeps every
member of \(H^\circ\) out of \(R_{\mathrm{ex}}\), because only a
counterexample excludes a valuation. Step 3 therefore always sees a
nonempty feasible set, so it does not fail. Since \(V\) never answers
\(\mathsf{unknown}\) and there is no cap, (2) ends the run on
\(\mathsf{proved}\). The returned \(h\) is feasible for step 3 at that
iteration, hence minimizes \(\mathit{dist}(\cdot, \theta_{P_1})\) over a
set that still contains \(H^\circ\). An attempt cap preserves (1) and (2)
and may drop (3).
\end{proof}

The bound \(|D|\) is a worst case. Side constraints shrink the search.
The guarantee covers only the holes of the given sketch. If no filling of
that sketch satisfies the contract, CEGIS fails, and a different structure
is a rewrite, not a closer \(h\).

\subsection{\textsc{BudgetNormalize}}

\begin{algo}[\textsc{BudgetNormalize}]
\label{alg:budgetnormalize}
Input: a finite set of CEGIS seeds, baseline \(P_1\), and a terminating
checker for \(\mathcal{B}\).
Output: a (possibly empty) subset of the seeds, and \(P_1\) unmarked as a
seed. \textsc{GeneratePool} is not skipped.
\begin{enumerate}
  \item Apply the checker for \(\mathcal{B}\) to each seed and to
  \(P_1\).
  \item Keep those seeds for which the checker passes. Drop the rest.
  \item If none remain, record that the seed set is empty.
  \item Return the kept set. Do not disable \textsc{GeneratePool}.
\end{enumerate}
\end{algo}

\begin{lemma}[\textsc{BudgetNormalize}]
\label{lem:budgetnormalize}
\begin{enumerate}
  \item \emph{Termination.} If the checker terminates on each of finitely
  many programs, so does the algorithm.
  \item \emph{Filter.} It adds no program.
  \item \emph{Kept seeds.} Every kept seed still carries whatever \(V\)
  proof it had, and passes \(\mathcal{B}\).
  \item \emph{Baseline.} The checker is applied to \(P_1\), but \(P_1\) is
  never added to the kept seeds, so a miss does not remove it from
  \eqref{eq:winner}.
  \item \emph{Empty seeds.} An empty kept set does not skip
  \textsc{GeneratePool}.
\end{enumerate}
\end{lemma}
\begin{proof}
One checker run per seed and one for \(P_1\) is finite if each run is,
which is (1). The algorithm only retains or drops members of its input,
which is (2). A kept seed is a subset of the input, so a prior
\(\mathsf{proved}\) answer remains, and step 2 is exactly the checker for
\(\mathcal{B}\), which is (3). Step 1 runs the checker on \(P_1\) and
step 2 never adds it to the kept set, which is (4). Step 4 returns after recording
emptiness and does not disable the next stage, which is (5).
\end{proof}

The stage does not repair a seed that misses \(\mathcal{B}\). A later
rewrite in \textsc{GeneratePool} may still produce a different program that
fits.

\subsection{\textsc{GeneratePool}}

Let \(R_{\mathrm{pool}}\) be a finite catalog of representation rewrites,
disjoint in purpose from the TheoryLower catalog though both are subsets of
the generator's \(R\).

\begin{algo}[\textsc{GeneratePool}]
\label{alg:generatepool}
Input: baseline \(P_1\), catalog \(R_{\mathrm{pool}}\), any seeds that
survived \textsc{BudgetNormalize}, verifier \(V\), and the checker for
\(\mathcal{B}\).
Output: a finite set \(\mathcal{C}\) of programs.
\begin{enumerate}
  \item Set \(\mathcal{C}\) to the surviving seeds (empty if none).
  \item For each rewrite \(r \in R_{\mathrm{pool}}\) that applies to
  \(P_1\), form the sketch \(r(P_1)\), run
  Algorithm~\ref{alg:cegis} on it, and if CEGIS returns a program that
  \(V\) proved and that passes \(\mathcal{B}\), add that program to
  \(\mathcal{C}\).
  \item Return \(\mathcal{C}\). Do not rank by \(\mathcal{O}\) or by an
  estimated cost.
\end{enumerate}
\end{algo}

\begin{lemma}[\textsc{GeneratePool}]
\label{lem:generatepool}
\begin{enumerate}
  \item \emph{Termination.} The algorithm terminates if each CEGIS run
  does and the checker for \(\mathcal{B}\) returns on each tried program.
  \item \emph{Starts from \(P_1\).} Step 2 runs even when the seed set is
  empty, and even when TheoryLower returned \(\mathsf{stuck}\).
  \item \emph{Admission.} Every program in the returned \(\mathcal{C}\)
  is in \(\mathcal{A}\). Surviving seeds are already in \(\mathcal{A}\).
  \item \emph{No ranking.} The algorithm does not consult
  \(\mathcal{O}\).
\end{enumerate}
\end{lemma}
\begin{proof}
\(R_{\mathrm{pool}}\) is finite. Each CEGIS run terminates by
Theorem~\ref{thm:cegis}(2). Each budget check is one run of a terminating
checker, which is (1). Step 2
iterates that catalog from \(P_1\) with no guard on the seed set or on
TheoryLower's tag, which is (2). A program added in step 2 was obtained
from \(P_1\) by a rewrite from the catalog, filled by CEGIS, proved by
\(V\), and passed \(\mathcal{B}\), hence lies in \(\mathcal{A}\). Surviving
seeds were proved by an earlier CEGIS run and passed
\textsc{BudgetNormalize}, which is (3). Step 3 forbids ranking, and no
other step reads \(\mathcal{O}\), which is (4).
\end{proof}

The forward pass can miss a combination that pays only after a later
rewrite. Exploring sequences of rewrites would find those combinations at
a cost exponential in sequence length. Ranking remains
\eqref{eq:winner}.

\subsection{Winner Selection}

\begin{algo}[Winner Selection]
\label{alg:winner}
Input: \(P_1\), the set \(\mathcal{C}\) from
Algorithm~\ref{alg:generatepool}, backend
\(\mathcal{L}_{\text{Target}}\), and objective \(\mathcal{O}\).
Output: source \(P^*\) and binary
\(\mathrm{Bin}^* = \mathcal{L}_{\text{Target}}(P^*)\).
\begin{enumerate}
  \item Let \(\mathcal{A}\) be \(\mathcal{C}\) (every member already
  proved and inside \(\mathcal{B}\)).
  \item If the contract names no \(\mathcal{O}\), return the first member
  of \(\mathcal{A}\), or \(P_1\) if \(\mathcal{A}\) is empty.
  \item Lower \(P_1\) and every member of \(\mathcal{A}\) with
  \(\mathcal{L}_{\text{Target}}\).
  \item Measure each binary on the named workload and hardware, and
  return the source whose binary minimizes \(\mathcal{O}\), preferring
  \(P_1\) on ties, as in~\eqref{eq:winner}. If \(\mathcal{A}\) is empty,
  return \(P_1\).
\end{enumerate}
\end{algo}

\begin{lemma}[Winner Selection]
\label{lem:winner}
\begin{enumerate}
  \item \emph{Termination.} If lowering and measurement return on each of
  finitely many programs, so does the algorithm. If there is no
  \(\mathcal{O}\), there is no measurement loop.
  \item It returns \(P_1\) or a member of \(\mathcal{A}\).
  \item Measurement is not admission: a program not proved by \(V\) is not
  ranked. With no \(\mathcal{O}\), the first admitted program is emitted.
\end{enumerate}
The remaining claims are Theorem~\ref{thm:restricted}.
\end{lemma}
\begin{proof}
\(\mathcal{A}\) is finite by Lemma~\ref{lem:generatepool}. With no
\(\mathcal{O}\), step 2 returns immediately, which is (1). Otherwise one
lowering and one measurement per member of
\(\mathcal{A} \cup \{P_1\}\) is finite if each call returns, which is (1).
The \(\arg\min\) set is
\(\{\mathcal{L}_{\text{Target}}(P) \mid P \in \mathcal{A} \cup \{P_1\}\}\),
or \(\{P_1\}\) when \(\mathcal{A}\) is empty, or the first of
\(\mathcal{A}\) when there is no \(\mathcal{O}\), so the returned source is
\(P_1\) or a member of \(\mathcal{A}\), which is (2). Steps 2 and 4 rank
only admitted programs, which is (3).
\end{proof}

\end{document}
